% Worked solutions for Listas/lista-r8-governance-activism.md (15 exercises).
% Compile: pdflatex twice; check: pdffonts (Type 1 only, uni = yes).
% Every number is asserted in lista-r8-governance-activism-solutions_check.py.
\documentclass[11pt,a4paper]{article}
\usepackage[utf8]{inputenc}
\usepackage[T1]{fontenc}
\usepackage{lmodern}   % Type 1 vector fonts (replaces PK bitmaps)
\usepackage{cmap}      % ToUnicode CMap -> copyable text and accents
\usepackage[brazilian,english]{babel}
\usepackage{amsmath,amssymb,amsfonts}
\usepackage[a4paper,margin=2.2cm]{geometry}
\usepackage{booktabs}
\usepackage{enumitem}
\usepackage{fancyhdr}
\usepackage{xcolor}
\usepackage{microtype}
\usepackage{tikz}

\pagestyle{fancy}\fancyhf{}
\lhead{\footnotesize Finan\c{c}as II --- Corporativas (MPE Insper)}
\rhead{\footnotesize List R8 --- Worked solutions}
\cfoot{\footnotesize\thepage}
\renewcommand{\headrulewidth}{0.3pt}
\setlength{\headheight}{14pt}
\setlist{topsep=2pt,itemsep=1pt}

\newcommand{\exsec}[1]{\section*{\color{blue!55!black}#1}\addcontentsline{toc}{section}{#1}}
\newcommand{\notation}{\par\noindent\textbf{Notation.}}
\newcommand{\step}[2]{\par\medskip\noindent\textbf{Step #1 --- #2.}~}
\newcommand{\intuition}{\par\medskip\noindent\textbf{Intuition.}~}
\newcommand{\bookans}{\par\medskip\noindent\textbf{Book answer.}~}
\newcommand{\refbib}{\par\smallskip\noindent\textbf{Ref:}~}
\newcommand{\asked}{\par\medskip\noindent\textbf{What is asked.}~}
\newcommand{\approach}{\par\smallskip\noindent\textbf{Approach.}~}
\newcommand{\answer}[1]{\par\smallskip\noindent\fbox{\textbf{Answer #1.}}~}
\newcommand{\pitfalls}{\par\medskip\noindent\textbf{Common mistakes.}}
\newcommand{\pt}{\foreignlanguage{brazilian}}
\newcommand{\dpp}{\Delta p}

\usepackage[hyphens]{url}
%% STATEMENT-PREAMBLE BEGIN (tools/statements_into_solutions.py)
\usepackage{longtable,array,calc}
\providecommand{\tightlist}{\setlength{\itemsep}{0pt}\setlength{\parskip}{0pt}}
\providecommand{\href}[2]{#2}
\providecommand{\real}[1]{#1}
\makeatletter\@ifundefined{c@none}{\newcounter{none}}{}\makeatother  % pandoc's \LTcaptype{none}
\newcommand{\stmtbegin}{\par\medskip\noindent\textbf{Statement.}\par\nopagebreak\begingroup\small\setlength{\tabcolsep}{3pt}\sloppy}
\newcommand{\stmtend}{\par\endgroup\nopagebreak\noindent\rule{\linewidth}{0.4pt}\par\medskip}
\DeclareUnicodeCharacter{2192}{\ensuremath{\rightarrow}}
\DeclareUnicodeCharacter{2212}{\ensuremath{-}}
\DeclareUnicodeCharacter{2022}{\textbullet{}}
\DeclareUnicodeCharacter{2248}{\ensuremath{\approx}}
\DeclareUnicodeCharacter{2032}{\ensuremath{^{\prime}}}
\DeclareUnicodeCharacter{2081}{\ensuremath{_{1}}}
\DeclareUnicodeCharacter{03B8}{\ensuremath{\theta}}
\DeclareUnicodeCharacter{03BB}{\ensuremath{\lambda}}
\DeclareUnicodeCharacter{03C8}{\ensuremath{\psi}}
\DeclareUnicodeCharacter{03B1}{\ensuremath{\alpha}}
\DeclareUnicodeCharacter{25A0}{\rule{0.7em}{0.7em}}
%% STATEMENT-PREAMBLE END
\begin{document}
\begin{center}{\Large\bfseries Exercise list R8 --- Corporate governance, control rights and investor activism}\\[4pt]{\large Worked solutions}\end{center}

\noindent Statements are in \url{Listas/lista-r8-governance-activism.md}; each exercise is cited by number and printed page and its full text is reproduced right before its solution. Short answers: \url{lista-r8-governance-activism-solutions-short.pdf}. Editions are the local copies: Tirole, \emph{The Theory of Corporate Finance} (2005); Amaro de Matos, \emph{Theoretical Foundations of Corporate Finance} (2001); Grinblatt--Titman, \emph{Financial Markets and Corporate Strategy}, 2nd ed. (2002); Ross = Ross, Westerfield, Jordan \& Lamb, \emph{Fundamentos de Administra\c{c}\~ao Financeira}, 9th ed. (2013). Throughout the Tirole exercises everyone is risk neutral, the rate of interest is $0$ and the entrepreneur is protected by limited liability; $\dpp\equiv p_H-p_L$.

\tableofcontents

% =============================================================================
\exsec{1. Management opposes a hostile offer --- Ross CQ 1.14, p.~23}
%% BEGIN-STATEMENT 1
\stmtbegin
\textbf{Source:} Ross, Cap. 1, Questão 14 (Revisão de conceitos e questões instigantes), p.~23 (PDF
47) · Answer in book: no

\textbf{14. Problemas de agência {[}OA4{]}} Suponhamos que você tenha ações de uma empresa. O preço
atual por ação é \$25. Outra empresa acaba de anunciar que deseja comprar a sua empresa e que pagará
\$35 por ação para adquirir todas as ações em circulação. A administração de sua empresa começa
imediatamente a se opor a essa oferta hostil. A administração está agindo em busca do melhor
interesse dos acionistas? Por quê?
\stmtend
%% END-STATEMENT

\asked Shares trade at \$25; a bidder offers \$35 for all shares and management immediately fights the offer. Is management acting in the shareholders' best interest, and why (or why not)?

\approach This is the manager--shareholder agency problem (Ross Cap.~1, ``problemas de ag\^encia'') seen through the market for corporate control. The test is whether resisting can raise what shareholders get; anything else serves the managers' private benefits of control (jobs, perks, independence).

\notation
\begin{itemize}
  \item $P_0=25$: pre-bid price per share; $P_b=35$: bid per share.
  \item $V^*$: value per share management claims the firm is worth on its own (unobserved by shareholders).
\end{itemize}

\step{1}{Size of the offer} The premium over the market price is
\begin{align}
  \frac{P_b}{P_0}-1 &= \frac{(35)}{(25)}-1 \label{r8e1-prem}\\
                    &= 0.40, \nonumber
\end{align}
a gain of $35-25=10$ per share that shareholders lose if the bid is defeated and the price returns to \$25.

\step{2}{When resisting is in the shareholders' interest} Resisting helps shareholders only if it leads to a better deal: (a) management credibly expects the stand-alone value to exceed the bid, $V^*>35$, i.e.\ more than $25\times1.40=35$ per share, information the market price does not yet reflect; or (b) resistance is a bargaining tactic to make this bidder raise the price or to attract competing bidders.

\step{3}{When it is not} An \emph{immediate} rejection, before any counter-offer is sought or any evidence on $V^*$ is presented, fits the agency explanation better: managers of a target usually lose their jobs and control benefits after a takeover, and those losses are not shared with shareholders. Their opposition then protects private benefits at a cost of \$10 per share to the owners.

\answer{} Not necessarily. Unless management can show the firm is worth more than \$35 per share or is using resistance to extract a higher bid, opposing a 40\% premium (\$10 per share) serves the managers' interest in keeping their jobs and control, not the shareholders' --- a classic agency problem.

\intuition The takeover market is a governance mechanism: a bid at a premium is the outsiders' estimate of what better management (or synergies) is worth. Managers who block it impose the cost of keeping control on the owners.

\pitfalls
\begin{itemize}
  \item Saying ``management must accept any premium'': resistance that triggers an auction can raise the price.
  \item Forgetting the counterfactual: if the bid fails, the price falls back to about \$25, not \$35.
\end{itemize}
\bookans No answer in the book.
\refbib{Ross Cap.~1, \S1.4 ``O problema de ag\^encia e o controle da empresa'' (agency relationship, managerial goals, control by shareholders, takeovers), pp.~12--16.}

% =============================================================================
\exsec{2. Ownership structure and institutional investors --- Ross CQ 1.15, p.~23}
%% BEGIN-STATEMENT 2
\stmtbegin
\textbf{Source:} Ross, Cap. 1, Questão 15 (Revisão de conceitos e questões instigantes), p.~23 (PDF
47) · Answer in book: no

\textbf{15. Problemas de agência e propriedade corporativa {[}OA4{]}} A propriedade corporativa
varia em todo o mundo. Historicamente, indivíduos têm tido a propriedade da maioria das ações das
corporações estadunidenses. Na Alemanha e no Japão, porém, os bancos, outras grandes instituições
financeiras e outras empresas possuem a maioria das ações das empresas abertas. Você acha que os
problemas de agência podem ser mais ou menos sérios na Alemanha e no Japão do que nos Estados
Unidos? Por quê? Nos últimos anos, grandes instituições financeiras, como os fundos mútuos e os
fundos de pensão, têm se tornado os proprietários dominantes de ações, e essas instituições estão
mais ativas nas questões corporativas. Quais são as implicações dessa tendência para os problemas de
agência e o controle corporativo?
\stmtend
%% END-STATEMENT

\asked (a) With banks and other firms owning most shares in Germany and Japan, are agency problems more or less serious there than in the US, where individuals historically held most shares? (b) What does the rise of active institutional owners (mutual and pension funds) imply for agency problems and corporate control?

\approach The severity of the manager--shareholder conflict depends on who has the incentive and the power to monitor. Dispersed owners free-ride on each other's monitoring; a large owner internalizes a large share of the gains (Grinblatt--Titman Result~18.2; Tirole \S9.2.1, ``concentrated ownership by a monitor with a sufficient stake improves the control of management'').

\step{1}{Free riding with dispersed owners} An owner of a fraction $\alpha$ of the shares who spends $c$ to discipline management gains only $\alpha\,\Delta V$. With $\alpha$ tiny, $\alpha\,\Delta V<c$ even when $\Delta V\gg c$: nobody monitors and managers have wide discretion (US model).

\step{2}{(a) Bank and cross-ownership} Banks and allied firms hold large blocks, often also lend to the firm and sit on its board. Their $\alpha$ is large, they have inside information, so the classical manager--shareholder agency problem is \emph{less} serious. But two other conflicts become \emph{more} serious: (i) controlling vs.\ minority shareholders (the controlling owner may extract private benefits; Ross Cap.~1 lists this conflict explicitly); (ii) the bank's interests as creditor (low risk, continued lending) differ from those of shareholders, and cross-holdings can entrench managers who protect each other.

\step{3}{(b) Institutional activism} Large funds hold blocks big enough to make monitoring pay and can coordinate votes (proxy fights, board pressure, CEO removals such as American Express in 1993 in Grinblatt--Titman). Implications: weaker manager--shareholder agency problems, more effective shareholder control, more pressure on boards and pay. Limits: fund managers are themselves agents of their clients, may be short-term oriented, and still capture only part of the gains they create.

\answer{(a)} Less serious for the manager--shareholder conflict (large owners monitor), but more serious for the conflict between controlling and minority shareholders and for creditor-oriented decisions.
\answer{(b)} Active institutions reduce the free-rider problem and strengthen shareholder control over management, at the cost of a new layer of agency (the funds' own managers).

\intuition Concentration trades one agency problem for another: it solves ``who watches the managers'' and creates ``who watches the watcher''.

\pitfalls
\begin{itemize}
  \item Answering ``less serious'' without naming the new conflict (majority vs.\ minority, bank as creditor).
  \item Treating institutions as owners with no agency problem of their own.
\end{itemize}
\bookans No answer in the book.
\refbib{Ross Cap.~1, \S1.4, pp.~12--16, and the chapter summary (conflicts between controlling and non-controlling shareholders), p.~22; Grinblatt--Titman Ch.~18, \S18.1 ``Can Financial Institutions Mitigate the Free-Rider Problem?'' and ``Changes in Corporate Governance'', pp.~632--633; Tirole \S9.2.1, p.~359.}

% =============================================================================
\exsec{3. Determinants of ownership concentration --- Grinblatt \& Titman 18.2, p.~653}
%% BEGIN-STATEMENT 3
\stmtbegin
\textbf{Source:} Grinblatt, Ch. 18, Exercise 18.2, p.~653 (PDF 666) · Answer in book: no

18.2. Discuss the factors that determine whether firms are likely to have large ownership
concentrations.
\stmtend
%% END-STATEMENT

\asked List and explain the forces that make a firm more or less likely to have a large shareholder.

\approach Ownership concentration is chosen by trading off the benefits of a large stake (better monitoring, control benefits, incentives of founders) against its private costs (undiversified risk, wealth needed), within the legal and regulatory environment (Results~18.2, 18.3, 18.5).

\step{1}{Benefits of a large stake} (i) \emph{Monitoring gains}: a large holder captures enough of the value created by disciplining management (Result~18.2); these gains are larger where management is hard to monitor and agency problems severe (Demsetz--Lehn: media firms) and in uncertain environments where discretion must be granted and watched (Result~18.5). (ii) \emph{Private benefits of control}: larger for firms with intangible assets (``more fun to own a controlling interest in a baseball team or a movie studio \ldots\ than a copper mine'', p.~636). (iii) \emph{Founders' incentives}: retaining shares commits the entrepreneur to effort and fetches a better IPO price (Result~18.3); selling may signal bad news, and capital-gains taxes discourage selling.

\step{2}{Costs of a large stake} The large holder bears firm-specific risk alone while sharing the gains with everyone. In Example~18.1 a \$60m investment buys 20\% of a firm worth $60/0.20=300$; adding \$60m of value is $60/300=20\%$ over five years,
\begin{align}
  (1.20)^{1/5}-1 &= 0.0371, \label{r8e3-ann}
\end{align}
under 4\% a year --- too little to compensate for holding a stock twice as volatile as the market. Larger and riskier firms therefore tend to have less concentrated ownership (a given stake costs more wealth and more risk).

\step{3}{Institutions} Weak legal protection of outside shareholders forces concentration (only a controlling owner is safe from expropriation); regulation that caps institutional stakes (Glass--Steagall, the 5\% limit on funds) keeps ownership dispersed in the US, while German and Japanese banks hold large blocks.

\answer{} Concentration is higher when monitoring gains and private control benefits are large (hard-to-monitor, uncertain, intangible-asset firms; founder-led firms), when the firm is small and less risky (cheaper to hold an undiversified block), and when legal protection of minorities is weak or institutions are allowed to hold blocks; it is lower when diversification costs dominate.

\intuition Every shareholder benefits from good management, but only the large holder pays for it in lost diversification; ownership is less concentrated than efficiency alone would dictate.

\pitfalls
\begin{itemize}
  \item Listing only benefits (monitoring) and forgetting the diversification cost that is borne privately.
  \item Ignoring the legal environment, which the chapter stresses across countries.
\end{itemize}
\bookans No answer in the book.
\refbib{Grinblatt--Titman Ch.~18, \S18.1 (Example~18.1, Result~18.2), pp.~631--634; \S18.2 (Result~18.3), pp.~634--636; Result~18.5, p.~639.}

% =============================================================================
\exsec{4. A law against blockholders' private benefits --- Grinblatt \& Titman 18.5, p.~653}
%% BEGIN-STATEMENT 4
\stmtbegin
\textbf{Source:} Grinblatt, Ch. 18, Exercise 18.5, p.~653 (PDF 666) · Answer in book: no

18.5. As a Washington policy analyst, you are asked to comment on a proposed law that would make it
more difficult for large outside shareholders to extract private benefits from the partial control
they can exert over management. How would such a law affect the incentives of outside shareholders
to monitor management?
\stmtend
%% END-STATEMENT

\asked How would a law that makes it harder for large outside shareholders to extract private benefits from partial control change their incentive to monitor management?

\approach The large shareholder bears the full cost of monitoring but receives only its pro-rata share of the value created (free-rider problem, Result~18.2). Private benefits of control are part of what makes monitoring worthwhile for him. Removing them lowers the return to monitoring.

\notation
\begin{itemize}
  \item $\alpha$: blockholder's fraction of shares; $\Delta V$: increase in firm value from monitoring; $c$: private cost of monitoring; $b_p$: private benefit he can extract.
\end{itemize}

\step{1}{Monitoring condition} The blockholder monitors iff
\begin{align}
  \alpha\,\Delta V + b_p &\ge c. \label{r8e4-ic}
\end{align}
Illustration: $\alpha=0.10$, $\Delta V=50$, $c=8$. Without private benefits, $\alpha\Delta V=(0.10)(50)=5<8$: no monitoring, although the other shareholders would gain $(0.90)(50)=45$. With $b_p=4$, $5+4=9\ge 8$: he monitors, and the others gain $45-4=41$ net of what he extracts. The smallest private benefit that sustains monitoring is $c-\alpha\Delta V=8-5=3$.

\step{2}{Effect of the law} The law lowers $b_p$. By \eqref{r8e4-ic}, monitoring falls (and fewer investors find it worth building a block). Minority shareholders gain because less is diverted per unit of monitoring, but lose the value monitoring created.

\step{3}{Net effect} The net effect is ambiguous: if private benefits are mostly ``compensation'' for monitoring (as in the example), the law reduces firm value; if they are mostly diversion with little monitoring behind them (Barclay--Holderness--Pontiff: closed-end funds with large holders trade at larger discounts), the law raises value. Private benefits also bias the blockholder's recommendations (compare Tirole's blockholder benefit $\xi$ in Exercise~12).

\answer{} It weakens outside blockholders' incentive to monitor: their pro-rata gain $\alpha\Delta V$ alone often falls short of the monitoring cost, and private benefits fill the gap. Minorities face less diversion but also less monitoring; the effect on firm value is ambiguous.

\intuition Private benefits are a ``wage'' for the public good of monitoring. Taxing the wage reduces the supply of the public good.

\pitfalls
\begin{itemize}
  \item Concluding the law is unambiguously good for minorities.
  \item Forgetting that the blockholder bears all of $c$ but gets only $\alpha$ of $\Delta V$.
\end{itemize}
\bookans No answer in the book.
\refbib{Grinblatt--Titman Ch.~18, \S18.1 (free-rider problem, Result~18.2), pp.~631--632; closed-end fund evidence, p.~636.}

% =============================================================================
\exsec{5. Board control and debt ratios --- Grinblatt \& Titman 18.4, p.~653}
%% BEGIN-STATEMENT 5
\stmtbegin
\textbf{Source:} Grinblatt, Ch. 18, Exercise 18.4, p.~653 (PDF 666) · Answer in book: no

18.4. Consider three similar firms that differ only in the extent to which they are controlled by
their boards of directors. In firm 1, the board has complete control of the investment decisions,
operating decisions, and financing choices. In firm 2, the board is unable to monitor investment and
operating decisions, but they do control financing decisions. In firm 3, the board has very little
control over either investment, operating, or financing decisions. Describe how debt ratios are
likely to differ in the three firms.
\stmtend
%% END-STATEMENT

\asked Rank the likely debt ratios of three otherwise similar firms: (1) board controls investment, operations and financing; (2) board controls only financing; (3) board controls almost nothing.

\approach Managers prefer less debt than shareholders because debt raises the risk of bankruptcy (job loss) and limits discretion; shareholders value debt as a disciplining device (Results~18.6--18.7, Jensen's free cash flow). Who sets financing, and whether other tools of control exist, determines leverage.

\step{1}{Firm 3 (management in control)} Management chooses financing in its own interest: low debt. Lowest debt ratio.

\step{2}{Firm 2 (board controls financing only)} The board cannot watch investments directly, so it uses leverage as the substitute: fixed payments drain free cash and the threat of bankruptcy deters value-destroying projects. Example~18.3 shows the mechanics: after a \$100m start, a second \$100m project is worth $250$, $175$ or $125$ (vs.\ $50$ without it). It is positive-NPV in the good and medium states ($175-50-100=25>0$) and negative in the bad state ($125-50-100=-25$). Senior debt $F$ with new money junior raises funds iff $V_{\text{with}}-F\ge 100$:
\begin{align}
  175-F \ge 100 &\iff F\le 75, \label{r8e5-up}\\
  125-F < 100 &\iff F> 25, \label{r8e5-lo}
\end{align}
so any $F\in(25,75]$ lets the firm invest only when it should. All-senior debt ($F=100$) leaves $175-100=75<100$ in the medium state and blocks a good project. Highest debt ratio.

\step{3}{Firm 1 (board controls everything)} The board can veto bad projects directly, so debt need not discipline; leverage is set by the usual trade-off (taxes vs.\ distress costs). Intermediate debt ratio, above firm 3's (managers' aversion to debt is overridden) and below firm 2's (no disciplining premium).

\answer{} Debt ratio: firm 2 $>$ firm 1 $>$ firm 3. Mehran's evidence (leverage rises with outside-shareholder influence and managerial ownership) supports Result~18.6.

\intuition Debt is a control right ``in disguise'': when the board cannot control investment, it controls the cash that would finance it.

\pitfalls
\begin{itemize}
  \item Ranking firm 1 highest because ``the board has most control'': with direct control, debt has less disciplining work to do.
  \item Pushing leverage too high: too much senior debt creates debt overhang and kills positive-NPV projects.
\end{itemize}
\bookans No answer in the book.
\refbib{Grinblatt--Titman Ch.~18, \S18.4 ``Capital Structure and Managerial Control'' (Results~18.6--18.8, Example~18.3), pp.~639--642.}

% =============================================================================
\exsec{6. Proxy contests: contenders vote for themselves --- Amaro de Matos 8.2.6 \#10, p.~223}
%% BEGIN-STATEMENT 6
\stmtbegin
\textbf{Source:} Amaro de Matos, Ch. 8, Exercise 8.2.6 \#10, p.~223 (PDF 241) · Answer in book: yes,
p.~278

\begin{enumerate}
\def\labelenumi{\arabic{enumi}.}
\setcounter{enumi}{9}
\tightlist
\item
  Show that the incumbent and the rival managements will always want to vote for themselves
  regardless of their private signals.
\end{enumerate}

\paragraph{Context needed to solve it}

\emph{From Amaro de Matos §8.2.4 ``The Choice of Takeover Methods'', pp.~214--216 (PDF 232--234):}

A paper by Harris and Raviv (1988) goes into greater depth in exploring the determinants of
corporate takeover methods and their outcomes and price effects.

In their model there are three types of agent: the incumbent management, a rival management, and a
large number of atomistic shareholders. To finance the \$1 investment required, the incumbent may
issue common stock and (nonvoting) debt. It is assumed that the only effect of debt is to increase
the fraction of votes controlled by one of the contenders. Since debt reduces the value of the
incumbent's control, in the absence of a takeover threat there will be no debt.

It is assumed that the management currently in control has one of two possible levels of ability
(say, 1 and 2), with probabilities \(k\) and \((1-k)\), respectively. The present value of cash
flows \(Y\) depends only on the ability of the management; its two possible values are denoted,
respectively, \(Y_1\) and \(Y_2\), with \(Y_1 > Y_2\).

The expected value of the firm is

\[Y_I = kY_1 + (1-k)Y_2.\]

The incumbent is supposed to have all its wealth \(W_I\) invested in the firm. Thus the fraction of
the firm owned by the incumbent is

\[\alpha = \frac{W_I^*}{Y_I} \quad\text{with } W_I^* = W_I + Y_I - 1.\]

What happens when there is a takeover threat? This model assumes that there are benefits to being in
control of the firm in addition to those that originate in net cash flows. In the event that there
is a rival management trying to take over, the incumbent, to the extent that he/she owns stock in
the firm, should not be opposed to the takeover if he/she believes that the rival has a higher
ability. Benefits from control, however, provide contenders with incentives to vote for themselves
regardless of their beliefs.

When the rival appears, the incumbent may choose, as an antitakeover device, to issue debt in
exchange for equity, with the entire exchange being completed before the rival purchases equity. The
value of the equity for a debt level \(D\) and controlling management of ability \(i = 1, 2\) is
denoted by \(E_i(D)\), and it is assumed that \(E_1 > E_2\) and \(E\) is decreasing with \(D\). Then
the incumbent's share of the equity becomes

\[\alpha_I(D) = \alpha Y(D)/E(D)\]

where \(Y(D)\) denotes the value of the firm when the face value of the debt is \(D\). Under Harris
and Raviv's assumptions given the ability of the incumbent management, the value of the firm is not
affected by the debt level. In other words, the dependence of \(Y\) on \(D\) appears only because
the debt level affects the determination of who controls the firm. It follows that
\(\alpha_I(0) = \alpha\).

Also, assuming that the rival management invests its wealth \(W_R\) in the firm's stock, and given
that the incumbent already has \(\alpha_I\), his/her share will be

\[\alpha_R(D) = \min\{W_R/E(D),\ 1 - \alpha_I(D)\},\]

and the passive investors will retain the rest:

\[\alpha_P(D) = 1 - \alpha_R(D) - \alpha_I(D).\]

The goal of the incumbent is to maximize the value of his/her shares plus the expected benefits of
control. Denoting by \(K(D)\) the benefits of control to the incumbent given a debt level \(D\), the
incumbent's problem is to choose \(D\) in such a way as to maximize

\[\alpha Y(D) + K(D). \qquad (8.16)\]

Given the resulting distribution of ownership \((\alpha_P, \alpha_R, \alpha_I)\), the incumbent's
share of the votes in any control contest consists of his/her own votes plus a fraction \(q_1\) or
\(q_2\) of the passive votes, depending on his/her ability \(i\), where \(q_1 > q_2\). Here, \(q_i\)
is the exogenous fraction of the passive shareholders voting for the incumbent, given his/her
ability is \(i\).

Hence, if the incumbent is of ability \(i\), he/she will have \(\alpha_I + q_i\alpha_P\) votes. It
is clear that this number depends on the debt level. But on the other hand, the outcome of the
takeover attempt depends on the number of votes on the side of the incumbent and therefore on the
debt level choice.

For instance, a \emph{successful tender offer} occurs in a situation in which the incumbent (even if
of higher ability) will not get 50\% of the votes. Formally, a tender offer is successful if \(D\)
is chosen so as to satisfy

\[\alpha_I(D) + q_2\alpha_P(D) \geq \frac{1}{2}. \qquad (8.17)\]

{[}\ldots{]} Finally, a proxy fight is a situation in which the incumbent gets more than 50\% of the
votes only if it is of type 1.
\stmtend
%% END-STATEMENT

\asked Show that in the Harris--Raviv model both the incumbent and the rival vote their own shares for themselves whatever their private information about who is the more able manager.

\approach In a control contest a contender's shares pay $\alpha_j E_i(D)$, with $i$ the ability of whoever controls; control also yields private benefits $K$. If the worst payoff in control beats the best payoff out of control, self-voting is (weakly) dominant whatever the beliefs. Harris and Raviv assume exactly what is needed (stated in the book's answer, p.~278): $E_2(D)+K_2(D)\ge E_1(D)$.

\notation
\begin{itemize}
  \item $j\in\{I,R\}$: incumbent, rival; $\alpha_j\in(0,1)$: fraction of the equity held by $j$ ($\alpha_I(D)$, $\alpha_R(D)$ in the model).
  \item $E_i(D)$: value of equity under a management of ability $i\in\{1,2\}$, $E_1>E_2$; $K_i(D)\ge 0$: control benefits of a manager of ability $i$, with $K_2$ the smaller one (worst case).
\end{itemize}

\step{1}{Payoffs} If $j$ wins control and has ability $i$, he gets $\alpha_jE_i(D)+K_i(D)$; if he loses, the other contender's ability $i'$ applies and he gets $\alpha_jE_{i'}(D)$. His private signal only changes his beliefs about $i$ and $i'$.

\step{2}{Worst case in control vs.\ best case out of control} The lowest in-control payoff is $\alpha_jE_2+K_2$; the highest out-of-control payoff is $\alpha_jE_1$. Their difference is
\begin{align}
  (\alpha_jE_2+K_2)-\alpha_jE_1 &= K_2-\alpha_j(E_1-E_2) \label{r8e6-diff}\\
  &\ge (E_1-E_2)-\alpha_j(E_1-E_2) \qquad\text{(using } K_2\ge E_1-E_2\text{)} \nonumber\\
  &= (1-\alpha_j)(E_1-E_2) > 0. \nonumber
\end{align}

\step{3}{Voting} A vote matters only when it is pivotal. In every pivotal event, \eqref{r8e6-diff} says winning beats losing for any realization of abilities, hence for any beliefs: voting for oneself is a dominant strategy for $I$ and for $R$. Numbers: $E_1=100$, $E_2=80$, $K_2=25$ ($80+25\ge100$), $\alpha=0.3$: worst in control $0.3(80)+25=49$ $>$ best out of control $0.3(100)=30$. If instead $K_2=5$, then $0.3(80)+5=29<30$: an incumbent who believed the rival more able would vote for the rival.

\answer{} Under $E_2+K_2\ge E_1$, $\alpha_jE_2+K_2-\alpha_jE_1\ge(1-\alpha_j)(E_1-E_2)>0$, so each contender prefers control in every state and votes for himself regardless of his signal.

\intuition Private benefits of control swamp the cash-flow stake: owning $\alpha_j<1$ of the equity, a contender internalizes only part of the value lost by an inferior management, but all of the control benefits. That is why the model treats the contenders' votes as fixed and lets the passive shareholders ($q_1$, $q_2$) decide.

\pitfalls
\begin{itemize}
  \item Comparing payoffs for the same ability in and out of control: the right comparison is worst-in vs.\ best-out, because the signal is private.
  \item Forgetting the factor $\alpha_j$: the needed condition is $K_2\ge\alpha_j(E_1-E_2)$, implied by the assumption since $\alpha_j\le1$.
\end{itemize}
\bookans Matches the book's answer (p.~278), which uses the same assumption and comparison.
\refbib{Amaro de Matos \S8.2.4 ``The Choice of Takeover Methods'', pp.~214--218; App.~C, p.~278.}

% =============================================================================
\exsec{7. The large shareholder: premium, cutoff and firm value --- Amaro de Matos 8.1.5 \#5, p.~205}
%% BEGIN-STATEMENT 7
\stmtbegin
\textbf{Source:} Amaro de Matos, Ch. 8, Exercise 8.1.5 \#5, p.~205 (PDF 223) · Answer in book: yes,
pp.~271--272

\begin{enumerate}
\def\labelenumi{\arabic{enumi}.}
\setcounter{enumi}{4}
\tightlist
\item
  When there is a large shareholder, explain and show what happens to the following variables as the
  large shareholder becomes larger:
\end{enumerate}

\begin{enumerate}
\def\labelenumi{(\alph{enumi})}
\item
  The tender offer, \(\pi(\alpha)\).
\item
  The cutoff value of \(Z\), \(Z^c(\alpha)\).
\item
  The market value of the firm, \(V(\alpha, q)\).
\end{enumerate}

Assume now that \(c = 0\) and let \(F(Z)\) be the uniform distribution on \([0, 1]\). Solve for the
values of \(\pi(\alpha)\) and the cutoff value \(Z^c(\alpha)\). Verify the relationships derived
earlier in the exercise.

\paragraph{Context needed to solve it}

\emph{From Amaro de Matos §8.1.2 ``The Role of a Large Shareholder'', pp.~196--198 (PDF 214--216):}

In the former model the bidder is an outsider who makes a profit at the expense of the shareholders,
who agree with the takeover in order to punish the existing management. It is known, however, that
large blockholders serve as an internal control on the managerial behavior and play an important
role in takeovers. This provides an alternative defensive mechanism against bad management, through
the possibility of proxy fights or tender offers from the large shareholders.

An extension of the former model seems to be in order, substituting for Assumption 8.1.1 with

\emph{Assumption 8.1.2.} There is one large minority shareholder who owns a given proportion of the
shares, and there is a fringe of atomistic shareholders.

This idea was formalized by Shleifer and Vishny (1986). Hereinafter the large shareholder will be
referred to as \(L\). Let \(\alpha\) be the proportion of the shares owned by \(L\). Since he/she is
a minority, \(0 < \alpha < \frac{1}{2}\). Also, let \(Z\) denote the random increase in the
discounted market value generated by the substitution of an inefficient management.

Assume that the increase in value comes from the fact that \(L\) has privileged access to a new
technology. Let \(I\) denote the probability of drawing a positive improvement value \(Z\) from a
probability distribution \(F(Z)\) with support on \((0, Z_{\max}]\). In this case, \(I\) can be
thought of as research intensity. This probability (or research intensity) comes at a cost \(c(I)\)
satisfying \(c' > 0\) and \(c'' > 0\). Under these conditions, the bidder will move to make an offer
if he/she can purchase \(\frac{1}{2} - \alpha\) of the shares. The expected market value under the
existing management being \(q\), the tender price must be higher than \(q\), say,

\[p = q + \pi,\]

where \(\pi\) must satisfy the condition

\[\frac{1}{2}Z - (0.5 - \alpha)\pi - c \geq 0, \qquad (8.2)\]

such that the bid is profitable for the raider. For the small shareholders the expected benefit from
a bid is

\[E[Z \mid 0.5Z - (0.5-\alpha)\pi - c \geq 0]\]

and therefore, to tender is the best strategy if

\[\pi - E[Z \mid Z - (1-2\alpha)\pi - 2c \geq 0] \geq 0. \qquad (8.3)\]

The role of \(L\) is to define the tender price by choosing the value of \(\pi\) that minimizes the
left-hand side of (8.3). The optimal value clearly depends on \(\alpha\). Let it be denoted by
\(\pi^*(\alpha)\). Notice that \(\pi^*(\alpha)\) is just the expected value of an improvement,
conditional on the takeover being profitable.

Shleifer and Vishny show that \(\pi^*(\alpha)\) is unique and is a decreasing function of
\(\alpha\). Notice that \(\alpha = 0\) replicates the case of diffuse ownership discussed already.
In this case, equation (8.2) implies that \(Z > \pi\) is a necessary condition for a takeover to
succeed. But under that condition no one tenders, reproducing the result of Grossman and Hart
(1990).

Defining \(Z^c(\alpha)\) as the cutoff value of \(Z\) that makes \(L\) indifferent between taking
over and not (that is to say, the smallest value of \(Z\) satisfying (8.2) with
\(\pi \equiv \pi^*(\alpha)\)), they show that \(Z^c(\alpha)\) is a strictly decreasing function.
Also, \(I^*(\alpha)\), the optimal probability of drawing a positive value of \(Z\), increases with
\(\alpha\), the proportion of shares held by \(L\). {[}\ldots{]}

To see how an increase in \(\alpha\) affects the market value of the firm, notice that this market
value may be written as

\[V(\alpha, q) = q + I^*(\alpha)\{1 - F[Z^c(\alpha)]\}E[Z \mid Z \geq Z^c(\alpha)].\]

This formula tells us that the value of the firm increases from the original value \(q\) by an
amount that is proportional to the product of three factors: (1) the optimal probability of drawing
a positive value of \(Z\), (2) the expected value of increase in discounted market value generated
by the substitution of an inefficient management, provided that it is above a certain minimum
threshold that makes the investment worthwhile, \(E[Z \mid Z \geq Z^c(\alpha)]\), and (3) the
probability that such an increase in profits is above that threshold, \(1 - F[Z^c(\alpha)]\).
{[}\ldots{]}

\emph{Proposition 8.1.1 (Shleifer and Vishny, 1986).} An increase in the proportion of shares held
by \(L\) results in a decrease in the takeover premium but an increase in the market value of the
firm.
\stmtend
%% END-STATEMENT

\asked (a)--(c) As the large shareholder's stake $\alpha$ rises, what happens to the premium $\pi^*(\alpha)$, the cutoff improvement $Z^c(\alpha)$ above which he bids, and firm value $V(\alpha,q)$? Then, with bid cost $c=0$ and $Z\sim U[0,1]$, solve for $\pi^*$ and $Z^c$ and verify.

\approach Shleifer--Vishny: the small shareholders tender only if the premium covers the expected improvement conditional on a bid; the blockholder free-rides less the more he owns, since he needs to buy only $\tfrac12-\alpha$ shares. Monotone comparative statics follow from (8.2)--(8.3); the uniform case gives closed forms.

\notation
\begin{itemize}
  \item $\alpha\in[0,\tfrac12)$: $L$'s stake; $Z\in(0,Z_{\max}]$: improvement, cdf $F$; $\pi$: premium over $q$; $c\ge0$: cost of the bid.
  \item $I$: research intensity (probability of drawing a $Z$), cost $c(I)$ with $c'>0$, $c''>0$.
  \item $h(\pi,\alpha)\equiv E[Z\mid Z\ge(1-2\alpha)\pi+2c]$: expected improvement given that $L$ bids.
\end{itemize}

\step{1}{(a) The premium} By (8.3) shareholders tender iff $\pi\ge h(\pi,\alpha)$, and $L$ picks the smallest such $\pi$:
\begin{align}
  \pi^*(\alpha) &= \min\{\pi:\ \pi\ge h(\pi,\alpha)\}. \label{r8e7-pistar}
\end{align}
For fixed $\pi$, the threshold $(1-2\alpha)\pi+2c$ falls as $\alpha$ rises, and the conditional mean of a positive variable above a lower threshold is lower: $h$ is decreasing in $\alpha$. So the set $\{\pi:\pi\ge h(\pi,\alpha)\}$ expands as $\alpha$ rises and its minimum cannot rise: $\pi^*$ is decreasing in $\alpha$.

\step{2}{(b) The cutoff} $L$ bids iff (8.2) holds at $\pi^*$, i.e.\ $Z\ge Z^c$ with
\begin{align}
  Z^c(\alpha) &= (1-2\alpha)\pi^*(\alpha)+2c. \label{r8e7-zc}
\end{align}
Both factors $(1-2\alpha)$ and $\pi^*(\alpha)$ fall with $\alpha$, so $Z^c$ falls: a larger blockholder bids for smaller improvements.

\step{3}{(c) Firm value} $V=q+I^*\int_{Z^c}^{Z_{\max}}Z\,dF(Z)$, since $\{1-F(Z^c)\}E[Z\mid Z\ge Z^c]=\int_{Z^c}Z\,dF$. A lower $Z^c$ adds positive improvements to the integral. Research: $L$'s expected gain per unit of $I$ is
\begin{align}
  G(\alpha) &= \int_{Z^c}^{Z_{\max}}\big[\tfrac12Z-(\tfrac12-\alpha)\pi^*-c\big]dF(Z) \nonumber\\
  &= \int_{Z^c}^{Z_{\max}}\tfrac12\,(Z-Z^c)\,dF(Z) \qquad\text{(by \eqref{r8e7-zc}, } (\tfrac12-\alpha)\pi^*+c=\tfrac12Z^c\text{)}, \label{r8e7-G}
\end{align}
with $dG/dZ^c=-\tfrac12[1-F(Z^c)]<0$. Since $Z^c$ falls with $\alpha$, $G$ rises; the first-order condition $c'(I^*)=G(\alpha)$ with $c''>0$ gives $I^*$ increasing. Both factors raise $V$: Proposition~8.1.1.

\step{4}{Uniform case, $c=0$} With $Z\sim U[0,1]$, $E[Z\mid Z\ge z]=(1+z)/2$. At the minimum in \eqref{r8e7-pistar} the condition binds, so $\pi^*=(1+Z^c)/2$, and \eqref{r8e7-zc} with $c=0$ gives $Z^c=(1-2\alpha)\pi^*$. Substituting the first into the second:
\begin{align}
  Z^c &= (1-2\alpha)\frac{1+Z^c}{2} \nonumber\\
  2Z^c &= (1-2\alpha)+(1-2\alpha)Z^c \nonumber\\
  Z^c\big[2-(1-2\alpha)\big] &= 1-2\alpha \nonumber\\
  Z^c(\alpha) &= \frac{1-2\alpha}{1+2\alpha}=\frac{0.5-\alpha}{0.5+\alpha}, \label{r8e7-zcu}\\
  \pi^*(\alpha) &= \frac12\Big(1+\frac{1-2\alpha}{1+2\alpha}\Big)=\frac12\cdot\frac{2}{1+2\alpha}=\frac{1}{1+2\alpha}. \label{r8e7-piu}
\end{align}
Check that it is the minimum: $h(\pi)=\tfrac12[1+(1-2\alpha)\pi]$, and $\pi\ge h(\pi)\iff\pi(1+2\alpha)/2\ge\tfrac12\iff\pi\ge1/(1+2\alpha)$.

\step{5}{Verify the comparative statics} Differentiating \eqref{r8e7-piu} and \eqref{r8e7-zcu}:
\begin{align}
  \frac{d\pi^*}{d\alpha} &= -\frac{2}{(1+2\alpha)^2}<0, \label{r8e7-dpi}\\
  \frac{dZ^c}{d\alpha} &= \frac{-2(1+2\alpha)-2(1-2\alpha)}{(1+2\alpha)^2}=-\frac{4}{(1+2\alpha)^2}=-\frac{1}{(0.5+\alpha)^2}<0. \label{r8e7-dzc}
\end{align}
Value: $\{1-F(Z^c)\}E[Z\mid Z\ge Z^c]=(1-Z^c)\frac{1+Z^c}{2}=\frac{1-(Z^c)^2}{2}$, increasing as $Z^c$ falls; and \eqref{r8e7-G} gives $G=\tfrac12\int_{Z^c}^1(Z-Z^c)dZ=\tfrac14(1-Z^c)^2$, also increasing. Numbers:
\begin{center}
\begin{tabular}{ccccc}
\toprule
$\alpha$ & $Z^c$ & $\pi^*$ & $(1-Z^{c2})/2$ & $G=\tfrac14(1-Z^c)^2$\\
\midrule
$0$ & $1$ & $1$ & $0$ & $0$\\
$0.1$ & $0.8/1.2=2/3$ & $1/1.2=0.8333$ & $5/18=0.2778$ & $1/36=0.0278$\\
$0.2$ & $0.6/1.4=3/7$ & $1/1.4=0.7143$ & $20/49=0.4082$ & $4/49=0.0816$\\
\bottomrule
\end{tabular}
\end{center}
With $c(I)=0.1\,I^2/2$, $I^*=G/0.1$: $I^*=0.2778$ and $0.8163$, so $V-q=I^*(1-Z^{c2})/2=0.0772$ at $\alpha=0.1$ and $0.3332$ at $\alpha=0.2$. At $\alpha=0$, $Z^c=1=Z_{\max}$: no takeover ever happens (Grossman--Hart).

\answer{(a)} $\pi^*(\alpha)$ decreases in $\alpha$; uniform case $\pi^*=1/(1+2\alpha)$, $d\pi^*/d\alpha=-2/(1+2\alpha)^2$.
\answer{(b)} $Z^c(\alpha)$ decreases; uniform case $Z^c=(0.5-\alpha)/(0.5+\alpha)$, $dZ^c/d\alpha=-1/(0.5+\alpha)^2$.
\answer{(c)} $V(\alpha,q)$ increases (more takeovers, larger expected improvement, more research $I^*$); uniform case $V=q+I^*(\alpha)\,[1-Z^c(\alpha)^2]/2$ with $I^*$ increasing.

\intuition The blockholder already owns the gains on $\alpha$ of the shares, so he needs to buy fewer shares at a premium and can afford to bid for smaller improvements. The small holders still demand the full expected improvement ($\pi^*=E[Z\mid\text{bid}]$), but that expectation falls as bids become more frequent.

\pitfalls
\begin{itemize}
  \item Treating $\pi$ as fixed: it is the conditional expectation of $Z$ given that a bid happens, and that conditioning set moves with $\alpha$.
  \item Mixing the bid cost $c$ in (8.2) with the research cost $c(I)$: ``$c=0$'' refers to the bid cost.
\end{itemize}
\bookans Matches the book's answer (pp.~271--272): $Z^c=(0.5-\alpha)/(0.5+\alpha)$, $\pi^*=1/(1+2\alpha)$, $\partial\pi^*/\partial\alpha=-2/(1+2\alpha)^2$, $\partial Z^c/\partial\alpha=-1/(0.5+\alpha)^2$. (The book's condition (C.7) prints ``$Z-(1-2\alpha)-2c$'', dropping the factor $\pi$; the solution uses $(1-2\alpha)\pi$ as in (8.3).) Parts (a)--(c) are referred to Shleifer--Vishny (1986); Steps 1--3 supply the argument.
\refbib{Amaro de Matos \S8.1.2 ``The Role of a Large Shareholder'', pp.~196--199; App.~C, pp.~271--272.}

% =============================================================================
\exsec{8. When investor control makes financing harder --- Tirole 10.5, p.~420}
%% BEGIN-STATEMENT 8
\stmtbegin
\textbf{Source:} Tirole, Ch. 10, Exercise 10.5, p.~420 (PDF 433) · Answer in book: yes, p.~613

\textbf{Exercise 10.5 (when investor control makes financing more difficult to secure).} The general
thrust of control rights theory is that investors are reassured, and so are more willing to lend, if
they have control rights over the firm. The purpose of this exercise is to build a counterexample in
which investor control is self-defeating and jeopardizes financing.

\begin{enumerate}
\def\labelenumi{(\roman{enumi})}
\tightlist
\item
  An entrepreneur has cash \(A\) and wants to invest \(I > A\) into a (fixed-size) project. The
  project yields \(R > 0\) with probability \(p\) and 0 with probability \(1-p\). The probability of
  success is \(p_H\) if the entrepreneur behaves and \(p_L = p_H - \Delta p\) (\(\Delta p > 0\)) if
  the entrepreneur misbehaves. The entrepreneur receives private benefit \(B > 0\) in the latter
  case, and 0 in the former case. All parties are risk neutral, the entrepreneur is protected by
  limited liability, and the rate of interest in the economy is 0.
\end{enumerate}

What is the necessary and sufficient condition for the entrepreneur to be able to obtain financing
from investors?

\begin{enumerate}
\def\labelenumi{(\roman{enumi})}
\setcounter{enumi}{1}
\tightlist
\item
  Now add a control right. This control right can raise the expected revenue in the case of
  misbehavior, but does nothing in the case of good behavior; namely, the holder of the control
  right can select an action (``damage control'') that raises the probability of success from
  \(p_L\) to \(p_L + \nu\) (\(\nu > 0\)) in the case of misbehavior, but keeps \(p_H\) constant.
  This interim action imposes a cost \(\gamma > 0\) on the entrepreneur. (If the action is not
  selected, the probabilities of success are as in question (i), and there is no private cost
  \(\gamma\).) The choice of action is simultaneous (say) with the entrepreneur's choice of effort.
\end{enumerate}

First assume ``entrepreneur control'' (the entrepreneur is given the right to select this action or
not). Write the two incentive constraints for the entrepreneur to behave. Show that, compared with
question (i), the pledgeable income remains the same if \(\nu B/(\Delta p) \leqslant \gamma\), and
is decreased otherwise.

\begin{enumerate}
\def\labelenumi{(\roman{enumi})}
\setcounter{enumi}{2}
\tightlist
\item
  Next consider ``investor control.'' Assume that when indifferent, the investors select the
  dominant strategy, i.e., the damage-control action (alternatively, one can assume that the action
  raises \(p_H\) as well, to \(p_H + \varepsilon\), where \(\varepsilon\) is arbitrarily small).
  Show that the financing condition is now
\end{enumerate}

\[p_H\left[R - \frac{B}{\Delta p - \nu}\right] \geqslant I - A.\]

Conclude that investor control, besides reducing NPV, may also make it more difficult for the
entrepreneur to secure financing.
\stmtend
%% END-STATEMENT

\asked (i) The financing condition in the basic fixed-investment model. (ii) With a ``damage control'' action under entrepreneur control: write the two incentive constraints and show that pledgeable income is unchanged if $\nu B/\dpp\le\gamma$ and lower otherwise. (iii) Under investor control, derive the financing condition $p_H[R-B/(\dpp-\nu)]\ge I-A$ and conclude.

\approach Pledgeable income is the expected return that can go to investors while keeping (IC$_b$). Control matters through the incentive constraint: here the action helps only after misbehavior, so whoever holds it changes the \emph{deviation} payoff, not the equilibrium payoff. Investor control, usually a way to raise pledgeable income (Aghion--Bolton), now insures the entrepreneur against her own shirking.

\notation
\begin{itemize}
  \item $R_b$: entrepreneur's reward in case of success; $\mathcal P$: pledgeable income; $\nu>0$: rise in success probability when damage control follows misbehavior; $\gamma>0$: its private cost to the entrepreneur.
  \item Assume $\nu<\dpp$ (otherwise no reward can deter shirking under investor control).
\end{itemize}

\step{1}{(i) Basic model} (IC$_b$): $p_HR_b\ge p_LR_b+B$, i.e.\ $R_b\ge B/\dpp$. Investors receive $p_H(R-R_b)$, maximal at $R_b=B/\dpp$:
\begin{align}
  \mathcal P_1 &= p_H\Big(R-\frac{B}{\dpp}\Big) \ge I-A. \label{r8e8-P1}
\end{align}

\step{2}{(ii) Entrepreneur control: the two constraints} If she behaves she never takes the action (it leaves $p_H$ unchanged and costs $\gamma$). If she misbehaves she may add it. Behaving must beat both deviations:
\begin{align}
  p_HR_b &\ge p_LR_b+B, \label{r8e8-ic1}\\
  p_HR_b &\ge (p_L+\nu)R_b+B-\gamma. \label{r8e8-ic2}
\end{align}
Rearranging \eqref{r8e8-ic2}: $(\dpp-\nu)R_b\ge B-\gamma$.

\step{3}{(ii) Does \eqref{r8e8-ic2} bind?} At the minimum reward of \eqref{r8e8-ic1}, $R_b=B/\dpp$, the left side of \eqref{r8e8-ic2} minus its right side is
\begin{align}
  p_H\frac{B}{\dpp}-(p_L+\nu)\frac{B}{\dpp}-B+\gamma &= \dpp\frac{B}{\dpp}-\nu\frac{B}{\dpp}-B+\gamma \nonumber\\
  &= \gamma-\nu\frac{B}{\dpp}. \label{r8e8-slack}
\end{align}
If $\nu B/\dpp\le\gamma$, \eqref{r8e8-ic2} holds at $R_b=B/\dpp$ and $\mathcal P=\mathcal P_1$. If $\nu B/\dpp>\gamma$, the reward must rise to $R_b=(B-\gamma)/(\dpp-\nu)$, which exceeds $B/\dpp$ because
\begin{align}
  \frac{B-\gamma}{\dpp-\nu}>\frac{B}{\dpp} \iff \dpp(B-\gamma)>B(\dpp-\nu) \iff \nu B>\gamma\dpp, \nonumber
\end{align}
so $\mathcal P=p_H[R-(B-\gamma)/(\dpp-\nu)]<\mathcal P_1$.

\step{4}{(iii) Investor control} Investors take damage control in every case (weakly dominant: it never lowers their income). The entrepreneur then bears $\gamma$ whether she behaves or not:
\begin{align}
  p_HR_b-\gamma &\ge (p_L+\nu)R_b+B-\gamma \label{r8e8-ic3}\\
  (\dpp-\nu)R_b &\ge B, \nonumber\\
  \mathcal P_2 &= p_H\Big(R-\frac{B}{\dpp-\nu}\Big)\ge I-A. \label{r8e8-P2}
\end{align}
Since $B/(\dpp-\nu)>B/\dpp$, $\mathcal P_2<\mathcal P_1$; and the NPV falls by $\gamma$ (the action is taken on the equilibrium path, where it does nothing).

\step{5}{Numbers} $p_H=0.8$, $p_L=0.5$ ($\dpp=0.3$), $B=0.6$, $R=5$, $\nu=0.1$, $I-A=2$. $B/\dpp=2$, $\mathcal P_1=0.8(5-2)=2.4$; $\nu B/\dpp=0.2$. With $\gamma=0.25\ge0.2$: $\mathcal P=2.4$. With $\gamma=0.15<0.2$: $R_b=(0.6-0.15)/(0.3-0.1)=2.25$, $\mathcal P=0.8(5-2.25)=2.2$. Investor control: $B/(\dpp-\nu)=0.6/0.2=3$, $\mathcal P_2=0.8(5-3)=1.6<2\le2.2$: financed under entrepreneur control, not under investor control.

\answer{(i)} $p_H(R-B/\dpp)\ge I-A$.
\answer{(ii)} Constraints \eqref{r8e8-ic1}--\eqref{r8e8-ic2}; $\mathcal P=\mathcal P_1$ if $\nu B/\dpp\le\gamma$, $\mathcal P=p_H[R-(B-\gamma)/(\dpp-\nu)]<\mathcal P_1$ otherwise.
\answer{(iii)} $p_H[R-B/(\dpp-\nu)]\ge I-A$; investor control lowers both NPV (by $\gamma$) and pledgeable income, so it can make financing impossible.

\intuition Damage control is insurance against shirking. When investors hold it they cannot commit not to use it, so misbehavior becomes less costly to the entrepreneur and she must be paid more to behave. Control rights help investors only when the action they will take does not soften the manager's punishment.

\pitfalls
\begin{itemize}
  \item Writing only one IC under entrepreneur control: the relevant deviation is shirking \emph{plus} damage control.
  \item Forgetting $\gamma$ on both sides of \eqref{r8e8-ic3}: under investor control it is borne regardless of effort and cancels.
\end{itemize}
\bookans Matches the book's answer (p.~613): $\mathcal P_1$; (2) $p_HR_b\ge(p_L+\nu)R_b+B-\gamma$; $\mathcal P_2=p_H[R-B/(\dpp-\nu)]$.
\refbib{Tirole \S10.2.1, pp.~389--391 (pledgeable income and control); Answers, p.~613.}

% =============================================================================
\exsec{9. Limited versus extended investor control --- Tirole 10.7, p.~421}
%% BEGIN-STATEMENT 9
\stmtbegin
\textbf{Source:} Tirole, Ch. 10, Exercise 10.7, p.~421 (PDF 434) · Answer in book: yes, pp.~613--614

\textbf{Exercise 10.7 (extent of control).} A simple variation on the basic model of Section 10.2.1
involves a choice between limited investor control and extended investor control, rather than
between entrepreneur control and investor control. Suppose, in the model of Section 10.2.1, that
entrepreneur control is out of the picture (after you finish the exercise, you may want to think
about a sufficient condition for this to be case), but that there are two degrees of investor
control:

\emph{Limited.} The action taken then increases the probability of success by \(\tau_A > 0\) and
inflicts cost \(\gamma_A > 0\) on insiders.

\emph{Extended} (investors have control over a wide set of actions). The selected action then
increases the probability of success by \(\tau_B > \tau_A\) and inflicts cost
\(\gamma_B > \gamma_A\) on insiders.

Assume that

\[\tau_A R - \gamma_A > \tau_B R - \gamma_B.\]

Find conditions under which limited or extended investor control prevails.

\paragraph{Context needed to solve it}

\emph{From Tirole §10.2.1 ``The Aghion--Bolton Model'', pp.~389--391 (PDF 402--404):}

To illustrate this in the simplest possible way, let us return to the basic (fixed-size) model of
Chapter 3: to finance her project, the entrepreneur must borrow the difference between the
investment cost \(I\) and her net worth \(A\). The project succeeds (and then yields \(R\)) with
probability \(p\), where \(p = p_H\) if she behaves and \(p_L\) if she misbehaves (and then takes
private benefit \(B\)); otherwise the project fails and yields nothing. Let us further introduce the
possibility of taking an interim action that

\begin{enumerate}
\def\labelenumi{(\roman{enumi})}
\item
  raises the probability of success uniformly by \(\tau > 0\) (so the probability of success becomes
  \(p_H + \tau\) or \(p_L + \tau\), depending on the entrepreneur's behavior, if the action is
  taken, and remains \(p_H\) or \(p_L\) if the status quo action is selected); and
\item
  engenders private cost \(\gamma > 0\) for the entrepreneur (or, more generally, the firm's
  insiders).
\end{enumerate}

{[}\ldots{]} We assume that this interim action cannot be contracted upon at the initial (financing)
stage. By contrast, the parties can contract on who is entitled to decide. {[}\ldots{]}

Figure 10.1 (timeline, left to right): \emph{Financing stage} --- project costs \(I\); entrepreneur
has equity \(A < I\); borrows \(I - A\). → \emph{Interim action} --- choice between status quo
action (probability of success is \(p\)), and profit-enhancing action (probability of success is
\(p + \tau\)). → \emph{Moral-hazard stage} --- entrepreneur's choice affects the probability of
success: \(p = p_H\) (no private benefit) or \(p_L\) (private benefit \(B\)). → \emph{Outcome stage}
--- verifiable profit: \(R\) with probability \(p\) (or \(p + \tau\)), 0 with probability \(1 - p\)
(or \(1 - p - \tau\)).

The assumption that the profit-enhancing action is orthogonal to managerial moral hazard, i.e.,
raises the probability of success uniformly, simplifies the analysis since it does not affect the
incentive compatibility condition: if the profit-enhancing action is to be taken, then the incentive
constraint becomes

\[[(p_H + \tau) - (p_L + \tau)]R_b \geqslant B \quad\text{or}\quad (p_H - p_L)R_b \geqslant B,\]

where \(R_b\) is the borrower's reward in the case of success. Note that \(\tau\) does not enter the
incentive constraint. {[}\ldots{]}

Let us first focus on the interesting case in which the profit-enhancing action reduces aggregate
welfare and is thus first-best suboptimal:

\[\tau R < \gamma.\]

Suppose that the control right is given to the investors. Because they share part of the profit and
bear none of the cost, they indeed select the profit-enhancing action, resulting in pledgeable
income:

\[(p_H + \tau)\left(R - \frac{B}{(p_H + \tau) - (p_L + \tau)}\right) = (p_H + \tau)\left(R - \frac{B}{\Delta p}\right).\]

Because investors do not earn supranormal profits in a competitive capital market (they break even),
the NPV is also the entrepreneur's welfare when raising funds. This NPV must account for both the
increase \(\tau\) in the probability of success and the entrepreneur's cost \(\gamma\) associated
with the interim action:

\[U_b = \text{NPV} = (p_H + \tau)R - I - \gamma.\]

{[}\ldots{]} There is, however, no renegotiation since the entrepreneur has no money to compensate
investors for the loss of value on their claims.

Suppose in contrast that the entrepreneur does not relinquish control. Because \(R_b \leqslant R\),
\(\tau R_b < \gamma\), and therefore the entrepreneur does not pick the profit-enhancing action. In
words, the entrepreneur bears the entire cost and gets only part of the benefits of the
profit-enhancing action. The pledgeable income is, as in Chapter 3,

\[p_H\left(R - \frac{B}{\Delta p}\right),\]

and the NPV (i.e., the entrepreneur's payoff) is

\[p_H R - I > (p_H + \tau)R - I - \gamma.\]

As expected, allocating control to investors reduces the NPV by \(\gamma - \tau R > 0\). But it
increases pledgeable income by \(\tau[R - (B/\Delta p)]\).

Suppose now that

\[p_H\left(R - \frac{B}{\Delta p}\right) < I - A < (p_H + \tau)\left(R - \frac{B}{\Delta p}\right).\]

Then the entrepreneur has insufficient cash on hand and can raise funds only by relinquishing the
control right to the investors. {[}\ldots{]}

Lastly, we note that the allocation of control is a trivial issue when investor control is
first-best optimal, that is, when

\[\tau R > \gamma.\]

In this case, investor control increases both the NPV from \(p_H R - I\) to
\((p_H + \tau)R - I - \gamma\) and the pledgeable income from \(p_H(R - B/\Delta p)\) to
\((p_H + \tau)(R - B/\Delta p)\).
\stmtend
%% END-STATEMENT

\asked Investors always get control, either limited (action $\tau_A,\gamma_A$) or extended ($\tau_B>\tau_A$, $\gamma_B>\gamma_A$), with $\tau_AR-\gamma_A>\tau_BR-\gamma_B$. When is each chosen? And when is entrepreneur control ruled out?

\approach Aghion--Bolton logic (\S10.2.1): the entrepreneur maximizes NPV (her utility, since investors break even) subject to the investors' breakeven. Control is a currency: more investor control raises pledgeable income but costs NPV. She hands over the least control that gets the project financed.

\notation
\begin{itemize}
  \item $j\in\{A,B\}$; $\mathcal P_j=(p_H+\tau_j)(R-B/\dpp)$: pledgeable income; $U_j=(p_H+\tau_j)R-I-\gamma_j$: NPV; $A_j\equiv I-\mathcal P_j$: minimum net worth.
\end{itemize}

\step{1}{Incentives are unaffected} The action raises success probability uniformly, so (IC$_b$) stays $\dpp R_b\ge B$ and investors get at most $(p_H+\tau_j)(R-B/\dpp)$. Investors in control always take the action (they share the gain, bear no cost).

\step{2}{Ranking} NPV: $U_A-U_B=(\tau_AR-\gamma_A)-(\tau_BR-\gamma_B)>0$. Pledgeable income: $\mathcal P_B-\mathcal P_A=(\tau_B-\tau_A)(R-B/\dpp)>0$, hence $A_B<A_A$.

\step{3}{Choice} The entrepreneur takes limited control whenever it is financeable; extended only when it is the sole route:
\begin{align}
  \text{limited:}&\quad (p_H+\tau_A)\Big(R-\frac{B}{\dpp}\Big)\ge I-A \iff A\ge A_A, \label{r8e9-lim}\\
  \text{extended:}&\quad (p_H+\tau_B)\Big(R-\frac{B}{\dpp}\Big)\ge I-A>(p_H+\tau_A)\Big(R-\frac{B}{\dpp}\Big) \iff A_B\le A<A_A, \label{r8e9-ext}
\end{align}
no financing if $A<A_B$; in each case the chosen NPV must be $\ge0$. (As in footnote~10 of \S10.2.1, randomizing between the two degrees could leave investors exactly at breakeven.)

\step{4}{Ruling out entrepreneur control} Entrepreneur control gives NPV $p_HR-I$ and income $p_H(R-B/\dpp)$. If $\tau_AR-\gamma_A\ge0$, limited investor control has (weakly) higher NPV \emph{and} strictly higher pledgeable income, so it dominates: a sufficient condition. (Entrepreneur control may lose even if $\tau_AR<\gamma_A$, because it may not attract funding.)

\step{5}{Numbers} $p_H=0.6$, $p_L=0.35$, $B=1$ ($B/\dpp=4$), $R=12$, $I=7$, $\tau_A=0.1$, $\gamma_A=0.5$, $\tau_B=0.2$, $\gamma_B=1.8$ ($\tau_AR-\gamma_A=0.7>0.6=\tau_BR-\gamma_B$). $\mathcal P_A=0.7(8)=5.6$, $\mathcal P_B=0.8(8)=6.4$, so $A_A=7-5.6=1.4$, $A_B=7-6.4=0.6$; $U_A=8.4-7-0.5=0.9$, $U_B=9.6-7-1.8=0.8$. Limited control for $A\ge1.4$, extended for $0.6\le A<1.4$, none below $0.6$. Entrepreneur control ($\mathcal P_E=4.8$, needs $A\ge2.2$, NPV $0.2$) is dominated since $\tau_AR-\gamma_A=0.7\ge0$.

\answer{} Limited investor control iff $(p_H+\tau_A)(R-B/\dpp)\ge I-A$; extended iff that fails but $(p_H+\tau_B)(R-B/\dpp)\ge I-A$; otherwise no financing. Sufficient for ruling out entrepreneur control: $\tau_AR\ge\gamma_A$.

\intuition Strong balance sheets buy autonomy: firms with more net worth give up fewer control rights (bond covenants), weaker ones give up more (bank loans, venture capital).

\pitfalls
\begin{itemize}
  \item Choosing the regime with higher pledgeable income whenever feasible: the entrepreneur maximizes NPV, and income only needs to cover $I-A$.
  \item Letting $\tau$ enter the incentive constraint: it raises $p_H$ and $p_L$ equally.
\end{itemize}
\bookans Matches the book's answer (pp.~613--614), including the sufficient condition $\tau_AR-\gamma_A\ge0$.
\refbib{Tirole \S10.2.1 ``The Aghion--Bolton Model'', pp.~389--391; Answers, pp.~613--614.}

% =============================================================================
\exsec{10. Control and incentives: complements or substitutes --- Tirole 10.6, pp.~420--421}
%% BEGIN-STATEMENT 10
\stmtbegin
\textbf{Source:} Tirole, Ch. 10, Exercise 10.6, pp.~420--421 (PDF 433) · Answer in book: yes, p.~613

\textbf{Exercise 10.6 (complementarity or substitutability between control and incentives).} This
exercise pursues the agenda set in Exercise 10.5 by considering various forms of complementarity and
substitutability between the exercise of control rights and managerial incentives. It therefore
relaxes the assumption of separability between the two.

\begin{enumerate}
\def\labelenumi{(\roman{enumi})}
\tightlist
\item
  An entrepreneur has cash \(A\) and wants to invest \(I > A\) into a (fixed-size) project. The
  project yields \(R > 0\) with probability \(p\) and 0 with probability \(1-p\). The probability of
  success is \(p_H\) if the entrepreneur behaves and \(p_L = p_H - \Delta p\) (\(\Delta p > 0\)) if
  the entrepreneur misbehaves. The entrepreneur receives private benefit \(B > 0\) in the latter
  case, and 0 in the former case. All parties are risk neutral, the entrepreneur is protected by
  limited liability, and the rate of interest in the economy is 0.
\end{enumerate}

What is the necessary and sufficient condition for the entrepreneur to be able to obtain financing
from investors?

\begin{enumerate}
\def\labelenumi{(\roman{enumi})}
\setcounter{enumi}{1}
\tightlist
\item
  Now consider the possibility that a profit-enhancing action be chosen. For reasons of simplicity
  (but not for the sake of realism!), assume that this action is chosen simultaneously with effort.
  This action raises the probability of success to
\end{enumerate}

\begin{itemize}
\tightlist
\item
  \(p_H + \tau_H\) if the entrepreneur behaves, and
\item
  \(p_L + \tau_L\) if the entrepreneur misbehaves.
\end{itemize}

The action is indeed profit enhancing (\(\tau_L, \tau_H > 0\)) and is

\begin{itemize}
\tightlist
\item
  complementary with effort if \(\Delta\tau \equiv \tau_H - \tau_L > 0\),
\item
  substitutable with effort if \(\Delta\tau < 0\).
\end{itemize}

The action further inflicts a disutility \(\gamma\) on the manager, where

\[\max(\tau_L, \tau_H)\cdot R < \gamma.\]

Lastly, assume that the high effort must be induced in order for financing to occur.

Write the pledgeable income under investor control and entrepreneur control. When does investor
control increase the pledgeable income (and therefore facilitate financing)?

\paragraph{Context needed to solve it}

\textbf{Part 0 --- the exercise whose agenda this one pursues.} \emph{From Tirole, Ch. 10, Exercise
10.5, p.~420 (PDF 433) · Answer in book: yes, p.~613 (it is also exercise 8 of this list):}

\textbf{Exercise 10.5 (when investor control makes financing more difficult to secure).} The general
thrust of control rights theory is that investors are reassured, and so are more willing to lend, if
they have control rights over the firm. The purpose of this exercise is to build a counterexample in
which investor control is self-defeating and jeopardizes financing.

\begin{enumerate}
\def\labelenumi{(\roman{enumi})}
\tightlist
\item
  An entrepreneur has cash \(A\) and wants to invest \(I > A\) into a (fixed-size) project. The
  project yields \(R > 0\) with probability \(p\) and 0 with probability \(1-p\). The probability of
  success is \(p_H\) if the entrepreneur behaves and \(p_L = p_H - \Delta p\) (\(\Delta p > 0\)) if
  the entrepreneur misbehaves. The entrepreneur receives private benefit \(B > 0\) in the latter
  case, and 0 in the former case. All parties are risk neutral, the entrepreneur is protected by
  limited liability, and the rate of interest in the economy is 0.
\end{enumerate}

What is the necessary and sufficient condition for the entrepreneur to be able to obtain financing
from investors?

\begin{enumerate}
\def\labelenumi{(\roman{enumi})}
\setcounter{enumi}{1}
\tightlist
\item
  Now add a control right. This control right can raise the expected revenue in the case of
  misbehavior, but does nothing in the case of good behavior; namely, the holder of the control
  right can select an action (``damage control'') that raises the probability of success from
  \(p_L\) to \(p_L + \nu\) (\(\nu > 0\)) in the case of misbehavior, but keeps \(p_H\) constant.
  This interim action imposes a cost \(\gamma > 0\) on the entrepreneur. (If the action is not
  selected, the probabilities of success are as in question (i), and there is no private cost
  \(\gamma\).) The choice of action is simultaneous (say) with the entrepreneur's choice of effort.
\end{enumerate}

First assume ``entrepreneur control'' (the entrepreneur is given the right to select this action or
not). Write the two incentive constraints for the entrepreneur to behave. Show that, compared with
question (i), the pledgeable income remains the same if \(\nu B/(\Delta p) \leqslant \gamma\), and
is decreased otherwise.

\begin{enumerate}
\def\labelenumi{(\roman{enumi})}
\setcounter{enumi}{2}
\tightlist
\item
  Next consider ``investor control.'' Assume that when indifferent, the investors select the
  dominant strategy, i.e., the damage-control action (alternatively, one can assume that the action
  raises \(p_H\) as well, to \(p_H + \varepsilon\), where \(\varepsilon\) is arbitrarily small).
  Show that the financing condition is now
\end{enumerate}

\[p_H\left[R - \frac{B}{\Delta p - \nu}\right] \geqslant I - A.\]

Conclude that investor control, besides reducing NPV, may also make it more difficult for the
entrepreneur to secure financing.
\stmtend
%% END-STATEMENT

\asked (i) Financing condition of the basic model. (ii) With a profit-enhancing action that raises success to $p_H+\tau_H$ (effort) or $p_L+\tau_L$ (shirking) at disutility $\gamma>\max(\tau_L,\tau_H)R$: pledgeable income under entrepreneur and investor control, and when investor control raises it.

\approach As in Exercise~8 (Tirole 10.5, Part~0 of the block), control affects pledgeable income through the incentive constraint. Here the action is not orthogonal to effort: its effect differs across effort levels by $\Delta\tau$, which adds to or subtracts from the incentive wedge $\dpp$.

\notation
\begin{itemize}
  \item $\Delta\tau\equiv\tau_H-\tau_L$; assume $\dpp+\Delta\tau>0$; $\mathcal P_{EC}$, $\mathcal P_{IC}$: pledgeable income under entrepreneur and investor control.
\end{itemize}

\step{1}{(i)} As in \eqref{r8e8-P1}: $p_H(R-B/\dpp)\ge I-A$.

\step{2}{(ii) Entrepreneur control} With $R_b\le R$, $\tau_HR_b\le\tau_HR<\gamma$, so with effort she does not take the action: $(p_H+\tau_H)R_b-\gamma<p_HR_b$. Shirking with the action is also worse than shirking without it: $(p_L+\tau_L)R_b+B-\gamma<p_LR_b+B$ because $\tau_LR_b<\gamma$. Only $\dpp R_b\ge B$ matters:
\begin{align}
  \mathcal P_{EC} &= p_H\Big(R-\frac{B}{\dpp}\Big). \label{r8e10-PEC}
\end{align}

\step{3}{(ii) Investor control} Investors always take the action (it raises success under either effort). Then
\begin{align}
  (p_H+\tau_H)R_b-\gamma &\ge (p_L+\tau_L)R_b+B-\gamma \nonumber\\
  (\dpp+\Delta\tau)R_b &\ge B, \nonumber\\
  \mathcal P_{IC} &= (p_H+\tau_H)\Big(R-\frac{B}{\dpp+\Delta\tau}\Big). \label{r8e10-PIC}
\end{align}

\step{4}{When does investor control help?} If $\Delta\tau\ge0$: $B/(\dpp+\Delta\tau)\le B/\dpp$ and $p_H+\tau_H>p_H$, so $\mathcal P_{IC}>\mathcal P_{EC}$ (given $R>B/\dpp$). If $\Delta\tau<0$, solve $\mathcal P_{IC}>\mathcal P_{EC}$ for $\Delta\tau$:
\begin{align}
  R-\frac{B}{\dpp+\Delta\tau} &> \frac{p_H}{p_H+\tau_H}\Big(R-\frac{B}{\dpp}\Big) \nonumber\\
  \frac{B}{\dpp+\Delta\tau} &< \frac{(p_H+\tau_H)R-p_HR+p_HB/\dpp}{p_H+\tau_H}=\frac{\tau_HR+p_HB/\dpp}{p_H+\tau_H} \nonumber\\
  \dpp+\Delta\tau &> \frac{B(p_H+\tau_H)}{\tau_HR+p_HB/\dpp}. \label{r8e10-thr}
\end{align}
The right side is below $\dpp$ iff $Bp_H+B\tau_H<\dpp\tau_HR+p_HB$, i.e.\ iff $B<\dpp R$ (true whenever financing is possible at all): investor control helps unless substitutability is strong enough, $\Delta\tau<\Delta\tau^*<0$.

\step{5}{Numbers} $p_H=0.6$, $p_L=0.35$, $B=1$, $R=10$, $\tau_H=0.05$, $\gamma=1.5$. $\mathcal P_{EC}=0.6(10-4)=3.6$. Threshold \eqref{r8e10-thr}: $\dpp+\Delta\tau>0.65/(0.5+2.4)=0.2241$, i.e.\ $\Delta\tau^*=-0.0259$. With $\Delta\tau=-0.05$ ($\tau_L=0.10$): $\mathcal P_{IC}=0.65(10-1/0.2)=3.25<3.6$. With $\Delta\tau=0$: $0.65(6)=3.9$; with $\Delta\tau=0.02$: $0.65(10-1/0.27)=4.0926$.

\answer{(i)} $p_H(R-B/\dpp)\ge I-A$.
\answer{(ii)} $\mathcal P_{EC}=p_H(R-B/\dpp)$, $\mathcal P_{IC}=(p_H+\tau_H)[R-B/(\dpp+\Delta\tau)]$. Investor control raises pledgeable income always if $\Delta\tau\ge0$ (complements or separability), and if $\Delta\tau<0$ only when $\dpp+\Delta\tau>B(p_H+\tau_H)/(\tau_HR+p_HB/\dpp)$.

\intuition An action that pays off mostly when the manager shirks ($\Delta\tau<0$) cushions shirking; investor control then weakens incentives (Exercise~8 is the extreme case $\tau_H=0$, $\tau_L=\nu$). When the action pays off mostly with effort, it sharpens incentives and investor control is doubly useful.

\pitfalls
\begin{itemize}
  \item Using $(p_H+\tau_H)-(p_L+\tau_L)=\dpp$: that is the separable case only.
  \item Forgetting that $\gamma$ cancels in the investor-control IC.
\end{itemize}
\bookans Matches the book's answer (p.~613); the explicit threshold \eqref{r8e10-thr} is derived here (the book only says the condition ``may fail if $\Delta\tau$ is sufficiently negative'').
\refbib{Tirole \S10.2.1, pp.~389--391; Answers, p.~613.}

% =============================================================================
\exsec{11. Public debt, bank debt and junk bonds --- Tirole 9.1, p.~379}
%% BEGIN-STATEMENT 11
\stmtbegin
\textbf{Source:} Tirole, Ch. 9, Exercise 9.1, p.~379 (PDF 392) · Answer in book: yes, pp.~606--607

\textbf{Exercise 9.1 (low-quality public debt versus bank debt).} Consider the model of Section
9.2.1, except that the project has a positive NPV even if the entrepreneur misbehaves.

As usual, the entrepreneur is risk neutral and protected by limited liability. She has assets \(A\)
and must finance an investment of fixed size \(I > A\). The project yields \(R\) in the case of
success and 0 in the case of failure. The probability of success is \(p_H\) if the entrepreneur
behaves (no private benefit) and \(p_L\) if she misbehaves (private benefit \(B\)). Investors are
risk neutral and demand a 0 rate of return.

Instead of assuming that the project has positive NPV only in the case of good behavior, suppose
that

\[p_H R > p_L R + B > I.\]

Suppose further that there is a competitive supply of monitors and abundant monitoring capital. At
private cost \(c\), a monitor can reduce the entrepreneur's private benefit of misbehavior from
\(B\) to \(b\). Assume that

\[p_H\frac{B - b}{\Delta p} > c > (\Delta p)R - p_H\frac{b}{\Delta p}\]

and

\[(\Delta p)R > c + B.\]

Show that there exist thresholds \(A_1 < A_2 < A_3\) such that

\begin{itemize}
\tightlist
\item
  if \(A \geqslant A_3\), the firm issues high-quality public debt (public debt that has a high
  probability of being repaid);
\item
  if \(A_3 > A \geqslant A_2\), the firm borrows from a monitor (and from uninformed investors);
\item
  if \(A_2 > A \geqslant A_1\), the firm issues junk bonds (public debt that has a low probability
  of being repaid);
\item
  if \(A_1 > A\), the firm does not invest.
\end{itemize}

\paragraph{Context needed to solve it}

\emph{From Tirole §9.2.1 ``Benefit of Activism'', pp.~356--358 (PDF 369--371):}

To model the collection of prospective information, we start from the fixed-investment model of
Section 3.2 and add a monitor who can intervene in order to reduce the scope for moral hazard. A
risk-neutral entrepreneur with wealth \(A\) has a project costing \(I > A\) and must therefore
borrow \(I - A\) from investors. The project yields \(R\) when it succeeds and 0 when it fails. The
probability of success is \(p_H\) if the entrepreneur works and \(p_L = p_H - \Delta p\) if she
shirks.

Table 9.1

{\def\LTcaptype{none} % do not increment counter
\begin{longtable}[]{@{}llll@{}}
\toprule\noalign{}
& good project & bad project & Bad project \\
\midrule\noalign{}
\endhead
\bottomrule\noalign{}
\endlastfoot
Pr(success) & \(p_H\) & \(p_L\) & \(p_L\) \\
Private benefit & 0 & \(b\) & \(B\) \\
\end{longtable}
}

\emph{9.2.1.1 No Monitoring.} In the absence of monitoring, shirking provides private benefit \(B\).
Letting \(R_b\) denote the entrepreneur's reward in the case of success (she receives nothing in the
case of failure as she is protected by limited liability), incentive compatibility requires that

\[(\Delta p)R_b \geqslant B. \qquad (9.1)\]

Funding requires that the pledgeable income exceed the investors' investment:

\[p_H\left(R - \frac{B}{\Delta p}\right) \geqslant I - A. \qquad (9.2)\]

If this condition is satisfied, and given that the investors break even, the entrepreneur's utility
is equal to the project's NPV:

\[U_b = p_H R - I. \qquad (9.3)\]

\emph{9.2.1.2 Monitoring (With Fixed Intensity).} Let us now formalize the idea that monitoring
reduces the extent of moral hazard. A straightforward way of doing so is to assume that a monitor
can reduce the private benefit that can be enjoyed by the entrepreneur by shirking from \(B\) to
\(b < B\). The monitor must, however, bear an unobservable private monitoring cost \(c > 0\) in
order to achieve this reduction in private benefit.

An interpretation of this monitoring structure is as described in Table 9.1. The manager will have
to choose among a number of ex ante identical projects. The manager privately learns the payoffs
attached to each project. There are three relevant projects: (1) the ``good project,'' which yields
no private benefit and has probability of success \(p_H\); (2) the low-private-benefit ``bad
project,'' which yields private benefit \(b\) and has probability of success \(p_L\); and (3) the
high-private-benefit ``Bad project,'' which yields private benefit \(B\) and has probability of
success \(p_L\).

The monitor moves first. If he incurs effort cost \(c\), he is able to identify the
high-private-benefit Bad project and thus to prevent the entrepreneur from selecting it. But he
still cannot tell the other two projects apart, and so the entrepreneur, who can condition her
choice of project on the existence or absence of monitoring, can still choose the
low-private-benefit bad project if she wishes to. The monitor learns nothing when he does not incur
the monitoring cost \(c\); then, because the projects are still indistinguishable by the investors,
the entrepreneur can choose any of the three projects, as in the absence of monitoring {[}\ldots{]}.

Let us assume that the entrepreneur ``hires'' a monitor and that the monitor's incentives induce him
to monitor. The entrepreneur's private benefit from shirking is then equal to \(b\), and so, if
\(R_b\) denotes the entrepreneur's reward in the case of success, the entrepreneur works if and only
if

\[(\Delta p)R_b \geqslant b. \qquad (9.4)\]

We can further assume that \((\Delta p)R_b < B\); for, if \(R_b \geqslant B/\Delta p\), the
entrepreneur is induced to work even in the absence of monitoring. Monitoring is then useless.

The monitor too must be provided with an incentive scheme. We maintain the assumption of universal
risk neutrality, and so there is no loss of generality in assuming that the monitor receives a
reward \(R_m\) in the case of success and 0 in the case of failure (because of limited liability).
When not incurring cost \(c\), the monitor is unable to prevent the entrepreneur from shirking, and
so the incentive for monitoring is provided by an \(R_m\) satisfying

\[(\Delta p)R_m \geqslant c. \qquad (9.5)\]

\emph{Abundance of monitoring capital assumption.} Let us first assume that monitoring capital is
``abundant'' or ``not scarce.'' This means that there is a large supply of monitors who are willing
to invest their capital in the monitoring activity as long as they are as well-off by doing so as
with any other investment. They are thus willing to contribute to the firm's investment at level
\(I_m\) such that

\[p_H R_m - c = I_m. \qquad (9.6)\]

The monitor then obtains no rent, and receives net payment \((p_H R_m - I_m)\) equal to his
monitoring cost \(c\). {[}\ldots{]}

Nonmonitoring or uninformed investors are willing to fund the project if and only if

\[p_H(R - R_b - R_m) \geqslant I - A - I_m. \qquad (9.7)\]

And so, using (9.4)--(9.6), the necessary and sufficient condition for the project to be funded is

\[p_H\left(R - \frac{b}{\Delta p}\right) \geqslant I - A + c. \qquad (9.8)\]

That is, monitoring reduces the agency cost from \(p_H B/\Delta p\) to \(p_H b/\Delta p\), but adds
monitoring cost \(c\). Using (9.5) and (9.6), the monitor's stake \(R_m\) can be chosen equal to
\(c/(\Delta p)\) and the monitor's investment contribution equal to

\[I_m = \frac{p_L c}{\Delta p}.\]

{[}\ldots{]} When does the entrepreneur benefit from having a monitor? Because all investors,
including the monitor, obtain no rent from their relationship with the firm, the entrepreneur's
utility is equal to the project's NPV under monitoring:

\[U_b = p_H R - I - c. \qquad (9.10)\]
\stmtend
%% END-STATEMENT

\asked Show that, under the stated inequalities, there are thresholds $A_1<A_2<A_3$ such that the firm issues high-quality public debt ($A\ge A_3$), borrows from a monitor ($A_2\le A<A_3$), issues junk bonds ($A_1\le A<A_2$), or does not invest ($A<A_1$).

\approach Active monitoring (\S9.2.1): a monitor at cost $c$ cuts the private benefit from $B$ to $b$, lowering the agency rent and raising pledgeable income. New here: the project is positive-NPV even with shirking, so ``junk'' finance that lets her shirk is a third option. Each option has a feasibility threshold in $A$ and a utility (= NPV, since all investors break even); she picks the best feasible one.

\notation
\begin{itemize}
  \item $U^3_b,U^2_b,U^1_b$: entrepreneur's utility with high-quality debt, monitoring, junk; $R_m$, $I_m$: monitor's stake and contribution.
\end{itemize}

\step{1}{High-quality public debt} She must behave unmonitored: $\dpp R_b\ge B$ and $p_H(R-R_b)\ge I-A$. With $R_b=B/\dpp$, feasible iff
\begin{align}
  A\ge A_3 &\equiv I-p_H\Big(R-\frac{B}{\dpp}\Big), \qquad U^3_b=p_HR-I. \label{r8e11-A3}
\end{align}

\step{2}{Monitored finance} With abundant monitoring capital, $R_m=c/\dpp$ (from (9.5)) and $I_m=p_HR_m-c$ (9.6). Uninformed investors' condition (9.7) with $R_b=b/\dpp$:
\begin{align}
  p_H\Big(R-\frac{b}{\dpp}-\frac{c}{\dpp}\Big) &\ge I-A-\Big(p_H\frac{c}{\dpp}-c\Big) \nonumber\\
  p_H\Big(R-\frac{b}{\dpp}\Big) &\ge I-A+c \nonumber\\
  A\ge A_2 &\equiv I+c-p_H\Big(R-\frac{b}{\dpp}\Big), \qquad U^2_b=p_HR-I-c. \label{r8e11-A2}
\end{align}

\step{3}{Junk bonds} She shirks ($\dpp R_b<B$), success probability $p_L$, and investors break even at $p_L(R-R_b)=I-A$. Her utility is $p_LR_b+B-A=p_LR_b+B-(I-p_L(R-R_b))$, i.e.
\begin{align}
  U^1_b &= p_LR+B-I. \nonumber
\end{align}
The largest repayment is $p_LR$ (at $R_b=0$; she can keep any extra cash), so
\begin{align}
  A\ge A_1 &\equiv I-p_LR. \label{r8e11-A1}
\end{align}

\step{4}{Ordering of thresholds} Subtracting,
\begin{align}
  A_3-A_2 &= p_H\frac{B}{\dpp}-p_H\frac{b}{\dpp}-c = p_H\frac{B-b}{\dpp}-c>0, \label{r8e11-32}\\
  A_2-A_1 &= c-p_HR+p_H\frac{b}{\dpp}+p_LR = c-\Big[\dpp R-p_H\frac{b}{\dpp}\Big]>0, \label{r8e11-21}
\end{align}
both by the first assumption.

\step{5}{Ordering of utilities} $U^3_b-U^2_b=c>0$, $U^2_b-U^1_b=\dpp R-c-B>0$ by the second assumption, and $U^1_b>0$ by $p_LR+B>I$. Since a cheaper form of finance is preferred whenever feasible, and feasibility sets are nested ($A_1<A_2<A_3$), the pattern in the statement follows; for $A<A_1$ nothing is feasible.

\step{6}{Numbers} $p_H=0.8$, $p_L=0.5$, $R=7$, $B=0.9$, $b=0.6$, $c=0.6$, $I=4$ (checks: $5.6>4.4>4$; $0.8>0.6>0.5$; $2.1>1.5$). $A_3=4-0.8(7-3)=0.8$, $A_2=4.6-0.8(7-2)=0.6$, $A_1=4-3.5=0.5$; $U^3_b=1.6$, $U^2_b=1.0$, $U^1_b=0.4$; $R_m=2$, $I_m=1.0=p_Lc/\dpp$.

\answer{} $A_1=I-p_LR<A_2=I+c-p_H(R-b/\dpp)<A_3=I-p_H(R-B/\dpp)$, with utilities $p_HR-I>p_HR-I-c>p_LR+B-I>0$; the firm uses the cheapest feasible source.

\intuition Strong balance sheets borrow on the market; intermediate ones pay for monitoring (bank debt); weak ones can only borrow by letting the manager take private benefits, which the market prices as low-quality (junk) debt.

\pitfalls
\begin{itemize}
  \item Forgetting $I_m$: the monitor co-invests $p_HR_m-c$, which is why the condition reads $I-A+c$.
  \item Writing junk utility as $p_LR-I$: the private benefit $B$ is part of her payoff.
\end{itemize}
\bookans Matches the book's answer (pp.~606--607): same $A_1,A_2,A_3$ and $U^3_b>U^2_b>U^1_b>0$.
\refbib{Tirole \S9.2.1 ``Benefit of Activism'', pp.~356--359; Answers, pp.~606--607.}

% =============================================================================
\exsec{12. A large minority blockholder --- Tirole 10.3, pp.~419--420}
%% BEGIN-STATEMENT 12
\stmtbegin
\textbf{Source:} Tirole, Ch. 10, Exercise 10.3, pp.~419--420 (PDF 432) · Answer in book: yes, p.~612

\textbf{Exercise 10.3 (large minority blockholding).} Consider the active monitor model (see Chapter
9). The firm yields \(R\) in the case of success and 0 in the case of failure. The entrepreneur,
large shareholder, and small shareholders have shares \(s_1\), \(s_2\), and \(s_3\), respectively,
where \(s_1 + s_2 + s_3 = 1\). (To complete the model's description, one can, as in Chapter 9,
assume that \(s_1 R \geqslant b/\Delta p\) and \(s_2 R \geqslant c/\Delta p\), using the notation of
this chapter.) The small shareholders have formal control (one share bears one voting right and
\(s_3 > \frac{1}{2}\)).

The project can be modified in a countable number of ways (\(k = 0, 1, \ldots\)). Option 0 consists
in ``not modifying the project'' (this option is known to everyone). Options 1 through \(\infty\) do
modify the project; all but two of them have disastrous consequences for all parties (so taking a
modification at random is dominated by the status quo option 0). The two relevant modifications are
such that one increases the probability of success by \(\tau > 0\) and the other reduces it by
\(\mu > 0\). One involves a private cost \(\gamma > 0\) or a private benefit \(-\gamma\) for the
entrepreneur, with \((\tau + \mu)s_1 R < \gamma\), and the other no such cost. Lastly, an action may
involve a private benefit \(\xi\) for the large blockholder (or one of his subsidiary). There are
three states of nature, as shown in Table 10.2. In each state of nature, the left-hand payoffs
correspond to the (uninformed) investor-friendly modification and the right-hand payoffs to the
(uninformed) investor-hostile modification.

Table 10.2 Probabilities: \(\beta\) (state 1); \((1-\beta)\kappa\) (state 2);
\((1-\beta)(1-\kappa)\) (state 3).

{\def\LTcaptype{none} % do not increment counter
\begin{longtable}[]{@{}
  >{\raggedright\arraybackslash}p{(\linewidth - 12\tabcolsep) * \real{0.1429}}
  >{\raggedright\arraybackslash}p{(\linewidth - 12\tabcolsep) * \real{0.1429}}
  >{\raggedright\arraybackslash}p{(\linewidth - 12\tabcolsep) * \real{0.1429}}
  >{\raggedright\arraybackslash}p{(\linewidth - 12\tabcolsep) * \real{0.1429}}
  >{\raggedright\arraybackslash}p{(\linewidth - 12\tabcolsep) * \real{0.1429}}
  >{\raggedright\arraybackslash}p{(\linewidth - 12\tabcolsep) * \real{0.1429}}
  >{\raggedright\arraybackslash}p{(\linewidth - 12\tabcolsep) * \real{0.1429}}@{}}
\toprule\noalign{}
\begin{minipage}[b]{\linewidth}\raggedright
\end{minipage} & \begin{minipage}[b]{\linewidth}\raggedright
State 1 (friendly)
\end{minipage} & \begin{minipage}[b]{\linewidth}\raggedright
State 1 (hostile)
\end{minipage} & \begin{minipage}[b]{\linewidth}\raggedright
State 2 (friendly)
\end{minipage} & \begin{minipage}[b]{\linewidth}\raggedright
State 2 (hostile)
\end{minipage} & \begin{minipage}[b]{\linewidth}\raggedright
State 3 (friendly)
\end{minipage} & \begin{minipage}[b]{\linewidth}\raggedright
State 3 (hostile)
\end{minipage} \\
\midrule\noalign{}
\endhead
\bottomrule\noalign{}
\endlastfoot
Impact on probability of success & \(\tau\) & \(-\mu\) & \(\tau\) & \(-\mu\) & \(\tau\) &
\(-\mu\) \\
Private cost for entrepreneur & 0 & \(-\gamma\) & 0 & \(\gamma\) & 0 & \(\gamma\) \\
Private benefit for large blockholder & 0 & 0 & 0 & 0 & 0 & \(\xi\) \\
\end{longtable}
}

The timing goes as follows:

\begin{enumerate}
\def\labelenumi{(\arabic{enumi})}
\item
  The entrepreneur learns the two relevant modifications and their impact on payoffs, and makes a
  proposal to shareholders.
\item
  The large blockholder learns the relevant modifications and their impact on payoffs, and either
  seconds the entrepreneur's recommendation for a modification or makes a counterproposal.
\item
  Majority shareholders decide between the status quo and the recommendation(s).
\end{enumerate}

\begin{enumerate}
\def\labelenumi{(\roman{enumi})}
\item
  Predict the outcome in each state of nature.
\item
  Add a fourth state of nature in which the entrepreneur and large shareholder see eye-to-eye and
  both prefer a value-decreasing action (say, the payoffs in state 4 are as in state 2, except that
  for the second relevant action, ``\(\gamma\)'' becomes ``\(-\gamma\)'' and the large shareholder
  receives \(\xi\)). What would you predict?
\end{enumerate}

\paragraph{Context needed to solve it}

\emph{From Tirole §9.2.1.2 ``Monitoring (With Fixed Intensity)'', pp.~357--358 (PDF 370--371) ---
the active monitor model the exercise refers to:}

A straightforward way of doing so is to assume that a monitor can reduce the private benefit that
can be enjoyed by the entrepreneur by shirking from \(B\) to \(b < B\). The monitor must, however,
bear an unobservable private monitoring cost \(c > 0\) in order to achieve this reduction in private
benefit. {[}\ldots{]} There are three relevant projects: (1) the ``good project,'' which yields no
private benefit and has probability of success \(p_H\); (2) the low-private-benefit ``bad project,''
which yields private benefit \(b\) and has probability of success \(p_L\); and (3) the
high-private-benefit ``Bad project,'' which yields private benefit \(B\) and has probability of
success \(p_L\).

{[}\ldots{]} The entrepreneur's private benefit from shirking is then equal to \(b\), and so, if
\(R_b\) denotes the entrepreneur's reward in the case of success, the entrepreneur works if and only
if

\[(\Delta p)R_b \geqslant b. \qquad (9.4)\]

{[}\ldots{]} the monitor receives a reward \(R_m\) in the case of success and 0 in the case of
failure (because of limited liability). When not incurring cost \(c\), the monitor is unable to
prevent the entrepreneur from shirking, and so the incentive for monitoring is provided by an
\(R_m\) satisfying

\[(\Delta p)R_m \geqslant c. \qquad (9.5)\]

\emph{Note (list editor, not from the book):} with shares \(s_1\) and \(s_2\) of a profit \(R\), the
rewards are \(R_b = s_1 R\) and \(R_m = s_2 R\); this is how the statement's conditions
\(s_1 R \geqslant b/\Delta p\) and \(s_2 R \geqslant c/\Delta p\) map into (9.4) and (9.5).
\stmtend
%% END-STATEMENT

\asked (i) Predict which modification is adopted in each of the three states, given the entrepreneur's proposal, the blockholder's seconding or counterproposal, and the uninformed majority's vote. (ii) Add a state where both prefer the value-decreasing action, and predict.

\approach A cheap-talk/voting game: the uninformed majority (control: $s_3>\tfrac12$) infers the state from agreement or disagreement between two informed parties with biases ($\gamma$ for the entrepreneur, $\xi$ for the blockholder), then picks the proposal with the higher expected value. The blockholder's stake $s_2$ aligns him with investors only if $\xi$ is small. Monitoring by a large shareholder is valuable as long as its private benefit does not bias it.

\notation
\begin{itemize}
  \item Friendly modification ($+\tau$), hostile ($-\mu$); payoffs relative to status quo. Entrepreneur: $s_1R\tau$ (friendly), $-s_1R\mu\mp\gamma$ (hostile). Blockholder: $s_2R\tau$; $-s_2R\mu$ (+$\xi$ in state 3). Majority: $s_3R\tau$, $-s_3R\mu$.
  \item State probabilities: $\beta$, $(1-\beta)\kappa$, $(1-\beta)(1-\kappa)$; $V_E$, $V_L$: majority's expected value (per unit $s_3R$) of following the entrepreneur's or the blockholder's proposal after a disagreement.
\end{itemize}

\step{1}{Preferences} Entrepreneur: in state 1 hostile beats friendly iff $\gamma-s_1R\mu>s_1R\tau$, i.e.\ $\gamma>(\tau+\mu)s_1R$ (assumed); so she ranks hostile $>$ friendly $>$ status quo. In states 2--3 hostile costs her $\gamma$: friendly $>$ status quo $>$ hostile. Blockholder: friendly in states 1--2; in state 3 hostile iff $\xi>(\tau+\mu)s_2R$.

\step{2}{Aligned blockholder, $\xi\le(\tau+\mu)s_2R$} He always backs the friendly action: he seconds in states 2--3 and counters in state 1. Disagreement reveals state 1 and the majority follows him. Outcome: friendly ($+\tau$) in all states.

\step{3}{Biased blockholder, $\xi>(\tau+\mu)s_2R$: truthful play} Entrepreneur proposes her favourite; agreement occurs in state 2 only, disagreement in states 1 and 3, where the entrepreneur's proposal is hostile in state 1 and friendly in state 3, the blockholder's the reverse. Hence
\begin{align}
  V_E &= -\beta\mu+(1-\beta)(1-\kappa)\tau, \label{r8e12-VE}\\
  V_L &= \beta\tau-(1-\beta)(1-\kappa)\mu, \label{r8e12-VL}
\end{align}
and the majority takes $\max\{V_E,V_L,0\}$ (dividing by the disagreement probability does not change the ranking).

\step{4}{Which equilibrium} (a) If $V_L$ is the max: outcome friendly in states 1--2, hostile in state 3. The entrepreneur is indifferent in state 1 (her proposal is overruled either way), so truthful play is an equilibrium. (b) If $V_E$ is the max: hostile in state 1, friendly in states 2--3. (c) If $0$ is the max, truthful play breaks down: after disagreement the status quo is chosen, so in state 1 the entrepreneur prefers to propose friendly ($s_1R\tau>0$), which the blockholder seconds, and in state 3 the blockholder prefers to second friendly ($s_2R\tau>0$). The outcome is friendly in \emph{all} states; disagreement is off the path, and the status quo there is optimal for any belief $q$ on state 1 with $\tau/(\tau+\mu)\le q\le\mu/(\tau+\mu)$, an interval that contains the prior posterior $\beta/[\beta+(1-\beta)(1-\kappa)]$ exactly when $V_E\le0$ and $V_L\le0$.

\step{5}{Numbers} $\tau=0.1$, $\mu=0.2$, $\kappa=0.5$. $\beta=0.6$: $V_E=-0.12+0.02=-0.10$, $V_L=0.06-0.04=0.02$: case (a). $\beta=0.3$: $V_E=-0.06+0.035=-0.025$, $V_L=0.03-0.07=-0.04$: case (c), friendly everywhere (posterior $0.3/0.65=0.46\in[1/3,2/3]$).

\step{6}{(ii) Collusive state 4} In state 4 the entrepreneur gets $\gamma$ and the blockholder $\xi$ from the hostile action, so both rank it first: she proposes it, he seconds. Agreement now pools state 2 (friendly) with state 4 (hostile). With probabilities $\pi_2,\pi_4$, the majority rubber-stamps an agreed proposal iff
\begin{align}
  \pi_2\tau-\pi_4\mu &\ge 0. \label{r8e12-agree}
\end{align}
If it holds, the hostile action passes in state 4 (collusion against the minority); if not, the majority rejects joint recommendations and loses $\tau$ in state 2 too. With $\pi_2=0.2$: $\pi_4=0.05$ gives $0.02-0.01=0.01>0$ (adopt), $\pi_4=0.15$ gives $-0.01<0$ (reject).

\answer{(i)} Aligned blockholder: friendly in all states. Biased blockholder: agreement (state 2) $\Rightarrow$ friendly; disagreement $\Rightarrow$ the majority follows whichever of $V_E$, $V_L$ in \eqref{r8e12-VE}--\eqref{r8e12-VL} is larger and positive; if both are negative, both informed parties prefer to agree on the friendly action and it is adopted in every state.
\answer{(ii)} Management and blockholder jointly push the value-decreasing action in state 4; agreement stops certifying good proposals, and the minority either accepts it (when $\pi_2\tau\ge\pi_4\mu$) or rejects all joint proposals.

\intuition A blockholder is valuable as a second, independent opinion. Private benefits ($\xi$) make his opinion biased, and collusion with management destroys its information value: the minority can no longer tell a good joint proposal from a self-dealing one.

\pitfalls
\begin{itemize}
  \item Reading $-\gamma$ in state 1 as a cost: it is a private \emph{benefit} for the entrepreneur.
  \item Forgetting that the majority conditions on agreement or disagreement, not on the state.
\end{itemize}
\bookans The book (p.~612) gives the truthful equilibrium and the majority's choice $\max\{-\beta\mu+\tau(1-\kappa),\ \beta\tau-\mu(1-\beta)(1-\kappa),\ 0\}$. Its first term omits the factor $(1-\beta)$ on $\tau$ (the probability of state 3 is $(1-\beta)(1-\kappa)$), compare \eqref{r8e12-VE}; and when the maximum is $0$ the truthful profile is not an equilibrium (Step 4c). Part (ii) is not answered in the book.
\refbib{Tirole \S9.2.1, pp.~356--359 (active monitoring); \S10.2 (control rights), pp.~389--393; Answers, p.~612.}

% =============================================================================
\exsec{13. Uncertain managerial horizon and control rights --- Tirole 10.8, pp.~421--422}
%% BEGIN-STATEMENT 13
\stmtbegin
\textbf{Source:} Tirole, Ch. 10, Exercise 10.8, pp.~421--422 (PDF 434) · Answer in book: yes, p.~614

\textbf{Exercise 10.8 (uncertain managerial horizon and control rights).} This exercise considers
the allocation of control between investors and management when the entrepreneur has an uncertain
horizon.

We consider the fixed-investment model. The investment cost is \(I\) and the entrepreneur has only
\(A < I\). The entrepreneur is risk neutral and protected by limited liability; the investors are
risk neutral and demand rate of return equal to 0. The profit is equal to \(R\) in the case of
success and is 0 in the case of failure. In the absence of profit-enhancing action, the probability
of success is \(p\); when the profit-enhancing action is taken this probability becomes
\(p + \tau\), where \(\tau > 0\), but the action imposes a nonmonetary cost on insiders, \(\gamma\),
where

\[\gamma > \tau R.\]

As usual, \(p = p_H\) if the entrepreneur behaves (no private benefit) and \(p = p_L\) if she
misbehaves (private benefit \(B\)).

The twist relative to Chapter 10 is that the entrepreneur may not be able to run the project to
completion: with probability \(\lambda\), she must quit the firm for exogenous reasons. She learns
this after the investment is sunk, but before the moral-hazard stage. If the entrepreneur quits
(which will have probability \(\lambda\)), a new and \emph{cashless} manager will be brought in.
This manager is also risk neutral and protected by limited liability and has the same private
benefit, probabilities of success, and payoff in the case of success as the entrepreneur.

Figure 10.10 summarizes the timing.

Figure 10.10 (timeline): \emph{Date 0} --- project costs \(I\); entrepreneur has \(A < I\). →
\emph{Realization of entrepreneur's horizon}: with probability \(1 - \lambda\) the entrepreneur
stays; with probability \(\lambda\) a new manager is appointed. → \emph{Date 1} (on both branches)
--- moral hazard; exercise of control right (profit-enhancing action?). → Success \(R\) with
probability \(p\), or \(p + \tau\), or failure (0).

Let \(x\) and \(y\) in \([0, 1]\) denote the probabilities that investors receive control when the
entrepreneur and the replacement manager are in charge, respectively. And assume that

\[(p_H + \tau)\frac{B}{\Delta p} \geqslant \gamma\]

(interpret this assumption), and that

\[\rho_1 \equiv p_H R > I > \rho_0^+ \equiv (p_H + \tau)\left(R - \frac{B}{\Delta p}\right).\]

\begin{enumerate}
\def\labelenumi{(\roman{enumi})}
\tightlist
\item
  Assuming that incentives must be provided for good behavior (by either the entrepreneur or the
  replacement manager), write down the following.
\end{enumerate}

\begin{itemize}
\tightlist
\item
  The entrepreneur's utility. (Hint: this utility is slightly different from the project's social
  value. Why?)
\item
  The pledgeable income and the breakeven condition.
\end{itemize}

\begin{enumerate}
\def\labelenumi{(\roman{enumi})}
\setcounter{enumi}{1}
\tightlist
\item
  Argue that \(y = 1\). Find the conditions under which the project is undertaken. (Warning. Two
  conditions must be fulfilled: investors must be willing to finance it, and the entrepreneur must
  be willing to go ahead with it.)
\end{enumerate}

\paragraph{Context needed to solve it}

\emph{From Tirole §10.2.1 ``The Aghion--Bolton Model'', pp.~389--390 (PDF 402--403) --- what
investor control and entrepreneur control do in ``Chapter 10'':}

{[}\ldots{]} if the profit-enhancing action is to be taken, then the incentive constraint becomes

\[[(p_H + \tau) - (p_L + \tau)]R_b \geqslant B \quad\text{or}\quad (p_H - p_L)R_b \geqslant B,\]

where \(R_b\) is the borrower's reward in the case of success. Note that \(\tau\) does not enter the
incentive constraint. {[}\ldots{]}

Suppose that the control right is given to the investors. Because they share part of the profit and
bear none of the cost, they indeed select the profit-enhancing action, resulting in pledgeable
income:

\[(p_H + \tau)\left(R - \frac{B}{(p_H + \tau) - (p_L + \tau)}\right) = (p_H + \tau)\left(R - \frac{B}{\Delta p}\right).\]

{[}\ldots{]} Suppose in contrast that the entrepreneur does not relinquish control. Because
\(R_b \leqslant R\), \(\tau R_b < \gamma\), and therefore the entrepreneur does not pick the
profit-enhancing action. In words, the entrepreneur bears the entire cost and gets only part of the
benefits of the profit-enhancing action. The pledgeable income is, as in Chapter 3,

\[p_H\left(R - \frac{B}{\Delta p}\right),\]

and the NPV (i.e., the entrepreneur's payoff) is

\[p_H R - I > (p_H + \tau)R - I - \gamma.\]

\emph{From Tirole §10.2.1, footnote 10, p.~390 (PDF 403) --- random allocation of control:}

Optimally, the entrepreneur would want to relinquish control stochastically under the set of
inequalities just stated: that is, to give control to investors with probability \(x\), and to
retain control with probability \(1 - x\), such that \((p_H + x\tau)(R - B/\Delta p) = I - A\).
\stmtend
%% END-STATEMENT

\asked Interpret $(p_H+\tau)B/\dpp\ge\gamma$. (i) Write the entrepreneur's utility (and why it differs from social value), the pledgeable income and the breakeven condition. (ii) Show $y=1$ and find when the project is undertaken (financed and wanted by the entrepreneur).

\approach Aghion--Bolton with two possible managers. Investor control takes the action ($+\tau$, cost $\gamma$ to the manager in place). The entrepreneur's utility is the NPV \emph{minus} the rent of the replacement manager, who is not party to the initial contract; so control over the replacement is cheap for her.

\notation
\begin{itemize}
  \item $\lambda$: probability she quits; $x$, $y$: probabilities of investor control with the entrepreneur and with the new manager; $\rho_1=p_HR$, $\rho_0=p_H(R-B/\dpp)$, $\rho_0^+=(p_H+\tau)(R-B/\dpp)$.
\end{itemize}

\step{1}{The assumption} Any manager must get $R_b\ge B/\dpp$. Under investor control the new manager's payoff is $(p_H+\tau)R_b-\gamma\ge(p_H+\tau)B/\dpp-\gamma$; the assumption says it is $\ge0$: the cashless replacement accepts the job even though investors will impose $\gamma$ on him.

\step{2}{(i) Pledgeable income and breakeven} With $R_b=B/\dpp$ for whichever manager runs the firm,
\begin{align}
  \mathcal P(x,y) &= (1-\lambda)(p_H+\tau x)\Big(R-\frac{B}{\dpp}\Big)+\lambda(p_H+\tau y)\Big(R-\frac{B}{\dpp}\Big)\ge I-A. \label{r8e13-P}
\end{align}

\step{3}{(i) Entrepreneur's utility} She is paid only if she stays: her gross payoff is $(1-\lambda)[(p_H+\tau x)R_b-\gamma x]$. Investors break even, $\mathcal P=I-A$ with her $R_b$ chosen accordingly, so subtracting her outlay $A$:
\begin{align}
  U_b &= (1-\lambda)\big[(p_H+\tau x)R-\gamma x\big]+\lambda(p_H+\tau y)\Big(R-\frac{B}{\dpp}\Big)-I. \label{r8e13-Ub}
\end{align}
Social value adds the new manager's rent:
\begin{align}
  S &= (1-\lambda)\big[(p_H+\tau x)R-\gamma x\big]+\lambda\big[(p_H+\tau y)R-\gamma y\big]-I, \nonumber\\
  S-U_b &= \lambda\Big[(p_H+\tau y)\frac{B}{\dpp}-\gamma y\Big]\ge0. \label{r8e13-gap}
\end{align}

\step{4}{(ii) $y=1$} From \eqref{r8e13-Ub} and \eqref{r8e13-P},
\begin{align}
  \frac{\partial U_b}{\partial y} &= \lambda\tau\Big(R-\frac{B}{\dpp}\Big)>0,\qquad \frac{\partial\mathcal P}{\partial y}=\lambda\tau\Big(R-\frac{B}{\dpp}\Big)>0: \nonumber
\end{align}
investor control over the replacement raises both her utility (the cost $\gamma$ falls on the new manager's rent) and pledgeable income. So $y=1$.

\step{5}{(ii) Choice of $x$} $\partial U_b/\partial x=(1-\lambda)(\tau R-\gamma)<0$ while $\partial\mathcal P/\partial x=(1-\lambda)\tau(R-B/\dpp)>0$: set $x$ as low as financing allows. With $y=1$, $\mathcal P(0,1)=(1-\lambda)\rho_0+\lambda\rho_0^+$.
\begin{itemize}
  \item If $(1-\lambda)\rho_0+\lambda\rho_0^+\ge I-A$: $x=0$, and she invests iff $U_b\ge0$, i.e.\ $(1-\lambda)\rho_1+\lambda\rho_0^+\ge I$.
  \item Otherwise $x$ solves
  \begin{align}
    (1-\lambda)\rho_0+\lambda\rho_0^++(1-\lambda)\tau x\Big(R-\frac{B}{\dpp}\Big) &= I-A, \label{r8e13-x}
  \end{align}
  feasible iff $x\le1$, i.e.\ $\rho_0^+\ge I-A$; and she invests iff $U_b\ge0$ at that $x$.
\end{itemize}
Since $\rho_0^+<I$, $U_b\ge0$ with $x=0$ requires $\lambda\le(\rho_1-I)/(\rho_1-\rho_0^+)$: a long expected horizon.

\step{6}{Numbers} $p_H=0.6$, $\dpp=0.25$, $B=0.5$ ($B/\dpp=2$), $R=10$, $\tau=0.1$, $\gamma=1.2$, $I=5.8$, $\lambda=0.2$. Assumptions: $\gamma=1.2>\tau R=1$; $(0.7)(2)=1.4\ge1.2$; $\rho_1=6>5.8>\rho_0^+=5.6$; $\rho_0=4.8$. $\mathcal P(0,1)=0.8(4.8)+0.2(5.6)=4.96$. With $A=1.5$ ($I-A=4.3$): $x=0$, $U_b=0.8(6)+1.12-5.8=0.12\ge0$; social value $0.16$, the gap $0.2(1.4-1.2)=0.04$ is the new manager's rent. With $A=0.5$ ($I-A=5.3$): $0.64x=5.3-4.96$, $x=0.53125$, $U_b=0.8(6-0.2\times0.53125)+1.12-5.8=0.035\ge0$: undertaken. Here $x=0$ works for $A\ge5.8-4.96=0.84$, and $\lambda\le(6-5.8)/(6-5.6)=0.5$.

\answer{(i)} $U_b$ in \eqref{r8e13-Ub}; it is social value minus the replacement manager's rent \eqref{r8e13-gap}. Pledgeable income and breakeven: \eqref{r8e13-P}.
\answer{(ii)} $y=1$. If $(1-\lambda)\rho_0+\lambda\rho_0^+\ge I-A$: $x=0$ and the project goes ahead iff $(1-\lambda)\rho_1+\lambda\rho_0^+\ge I$. Otherwise $x\in(0,1]$ from \eqref{r8e13-x}, which requires $\rho_0^+\ge I-A$, and the project goes ahead iff $U_b\ge0$ at that $x$.

\intuition The entrepreneur keeps control for herself (she bears $\gamma$) but happily gives investors control over her successor, whose disutility she does not bear and who gets a rent anyway. This fits founders who keep control while in office but leave strong board rights for the succession.

\pitfalls
\begin{itemize}
  \item Equating $U_b$ with NPV: the replacement manager's rent is a leakage from her point of view.
  \item Forgetting the second condition: financing feasible does not mean $U_b\ge0$.
\end{itemize}
\bookans Matches the book's answer (p.~614). The book's equation for $x$ prints $(1-\lambda)\rho_0+\lambda\rho_0^++\tau x(R-B/\dpp)=I-A$, dropping the factor $(1-\lambda)$ on the $x$ term; from \eqref{r8e13-P} the factor belongs there, as in \eqref{r8e13-x}.
\refbib{Tirole \S10.2.1, pp.~389--391 (and footnote 10 on random control); Answers, p.~614.}

% =============================================================================
\exsec{14. A continuum of control rights --- Tirole 10.9, p.~422}
%% BEGIN-STATEMENT 14
\stmtbegin
\textbf{Source:} Tirole, Ch. 10, Exercise 10.9, p.~422 (PDF 435) · Answer in book: yes, pp.~614--615

\textbf{Exercise 10.9 (continuum of control rights).} This exercise extends the analysis of Section
10.2.2 to a continuum of control rights. As in Section 10.2.2, consider a risk-neutral entrepreneur
protected by limited liability. The entrepreneur has cash on hand \(A\) and wants to finance a
project with cost \(I > A\). The project yields \(R\) if it succeeds and 0 if it fails. Investors
are risk neutral and demand a rate of return equal to 0. There is a continuum of control rights,
where the decision attached to a control right can be thought of as a modification relative to the
initial project and is characterized by the pair \((t, g)\): \(t \gtrless 0\) is the increase in the
probability of success and \(g \gtrless 0\) is the private cost borne by the entrepreneur if the
decision is taken (the modification is made). Let \(F(t, g)\) denote the continuous joint
distribution over the space of control rights and \(E_F[\cdot]\) the expectations with respect to
distribution \(F\).

The probability of success is

\[p + \tau \equiv p + E_F[t\,x(t, g)],\]

where \(x(t, g) = 1\) if the decision \((t, g)\) is taken and 0 otherwise. Similarly, let

\[\gamma \equiv E_F[g\,x(t, g)].\]

Moral hazard is modeled in the usual way: \(p = p_H\) if the entrepreneur behaves (no private
benefit) and \(p = p_L\) if the entrepreneur misbehaves (and receives private benefit \(B\)). Assume
that the project can be funded only if the entrepreneur is provided with the incentive to behave.

\begin{enumerate}
\def\labelenumi{(\roman{enumi})}
\item
  Solve for the optimal policy \(x(\cdot, \cdot)\), assuming that the investors' breakeven
  constraint is binding (which it is for \(A\) small enough or \(I\) large enough).
\item
  Show that, as \(A\) decreases, \(\tau\) and \(\gamma\) increase.
\item
  Discuss the implementation of the optimal \(x(\cdot, \cdot)\) function.
\item
  Consider the degenerate case in which \(g\) is the same for all control rights (\(g > 0\)). Show
  that
\end{enumerate}

\[\frac{d^2\gamma}{d\tau^2} > 0.\]

\paragraph{Context needed to solve it}

\emph{From Tirole §10.2.2 ``Multiple Control Rights'', pp.~392--393 (PDF 405--406):}

In practice, there are multiple control rights to be divided between insiders and outsiders: product
design, day-to-day management, long-term strategic decisions, hiring decisions, mergers, alliance
building, etc. The analysis above is straightforwardly generalized. The intuition derived from the
single-control-right analysis suggests, and this section confirms, that, in the presence of multiple
control rights, it is always optimal for the entrepreneur to abandon all rights for which investor
control is first-best optimal as well as, possibly, some rights for which it is not. Again, the
optimal split of rights accounts not only for the value (NPV) impact of the allocation, but also for
its impact on pledgeable income.

More formally, and generalizing the framework of Section 10.2.1 to multiple rights, suppose there
are \(K\) dimensions of decision making and therefore \(K\) rights to allocate. Each right
\(k \in \{1, \ldots, K\}\) is characterized by the uniform increase \(\tau_k > 0\) in the
probability of success and the cost \(\gamma_k > 0\) borne by insiders if the course of action is
altered in dimension \(k\). The governance structure is now defined by the allocation
\(x \equiv \{x_1, \ldots, x_K\}\) of formal control rights, where \(x_k = 1\) if investors obtain
the control of decision \(k\) and \(x_k = 0\) if the entrepreneur retains control.

Let us maximize the NPV (that is, the entrepreneur's utility) subject to being able to secure
financing and to the entrepreneur's incentive compatibility constraint:

\[\max_{\{R_b, x\}}\left[\left(p_H + \sum_k \tau_k x_k\right)R - I - \sum_k \gamma_k x_k\right]\]

\[\text{s.t.}\quad \left(p_H + \sum_k \tau_k x_k\right)[R - R_b] \geqslant I - A, \qquad R_b \geqslant \frac{B}{\Delta p}.\]

The solution to this program has the following features in the interesting case in which the
financing constraint is binding: there exists a threshold \(\theta < 1\) such that investors receive
control over decision \(k\) if and only if their relative willingness to pay for that right exceeds
the threshold, that is, if and only if

\[\frac{\tau_k R}{\gamma_k} \geqslant \theta.\]

\emph{Footnote 15, p.~392 (PDF 405):} Let \(\mu\) denote the shadow price of the financing
constraint. In the interesting case in which \(\mu\) is strictly positive, then
\(R_b = B/\Delta p\), and the derivative of the Lagrangian with respect to \(x_k\) is equal to
\([\tau_k R - \gamma_k] + \mu\tau_k[R - B/\Delta p]\). The first term is just the (first-best)
efficiency of taking the profitability-enhancing stance on decision \(k\), while the second term
reflects the benefits in terms of relaxing the financing constraint. {[}\ldots{]}

\emph{Footnote 17, p.~392 (PDF 405):} {[}\ldots{]} The only rights for which the allocation is not
first-best efficient (i.e., those with \(\theta \leqslant \tau_k R/\gamma_k < 1\)) are rights that
are given to investors. Let us first assume that the entrepreneur has committed her cash on hand
\(A\) and borrowed only what is needed to fund the investment cost, \(I - A\) (it turns out to be
optimal to do so). Then the entrepreneur no longer has any means of payment after securing
financing, the entrepreneur has no way to compensate the investors for relinquishing these rights.
In other words, there are potential gains from trade ex post, but these gains cannot be reaped in
the absence of compensating transfer (technically, utility is ``nontransferable''). {[}\ldots{]} We
conclude that the optimal allocation of rights obtained above is not renegotiated.

\emph{From p.~393 (PDF 406):} The entrepreneur's incentive constraint is then binding
(\(R_b = B/\Delta p\)). Intuitively, the entrepreneur prefers to pay investors in the ``efficient''
currency (cash, which is one-for-one) than in the ``inefficient'' one (control rights, whose
transfer involves a deadweight loss). {[}\ldots{]} An important implication of this analysis is
that, ceteris paribus, firms with stronger balance sheets (say, with a higher \(A\); see Section
3.2.2 for a broader definition of balance-sheet strength) abandon fewer rights.
\stmtend
%% END-STATEMENT

\asked (i) Optimal rule $x(t,g)$ with binding breakeven. (ii) Show $\tau$ and $\gamma$ rise as $A$ falls. (iii) How to implement $x$ by allocating rights. (iv) With common $g>0$, show $d^2\gamma/d\tau^2>0$.

\approach Section~10.2.2's program with a continuum of rights: maximize NPV subject to breakeven and (IC$_b$). The Lagrangian is linear in each $x(t,g)$, so the solution is bang-bang with a cutoff line through the origin whose slope rises with the shadow price of funds.

\notation
\begin{itemize}
  \item $\mu\ge0$: multiplier on the breakeven constraint; $k(\mu)\equiv R+\mu(R-B/\dpp)$; $H$, $h$: cdf and density of $t$ in (iv).
\end{itemize}

\step{1}{(i) $R_b=B/\dpp$} $R_b$ does not enter the NPV and lowering it relaxes the breakeven; with a binding breakeven, (IC$_b$) binds.

\step{2}{(i) Lagrangian} With $R_b=B/\dpp$,
\begin{align}
  \mathcal L &= \big[p_H+E_F(tx)\big]R-I-E_F(gx)+\mu\Big\{\big[p_H+E_F(tx)\big]\Big(R-\frac{B}{\dpp}\Big)-(I-A)\Big\}. \nonumber
\end{align}
The derivative with respect to $x(t,g)$ at each point is
\begin{align}
  \frac{\partial\mathcal L}{\partial x(t,g)} &= tR-g+\mu t\Big(R-\frac{B}{\dpp}\Big) = t\,k(\mu)-g, \nonumber
\end{align}
so
\begin{align}
  x(t,g)=1 &\iff g\le k(\mu)\,t. \label{r8e14-rule}
\end{align}
In the $(t,g)$ plane: take every modification below the line $g=k(\mu)t$ through the origin, slope $k(\mu)\ge R$; $\mu$ is pinned down by $[p_H+\tau(\mu)](R-B/\dpp)=I-A$. First best is $\mu=0$, slope $R$.

\step{3}{(ii) Comparative statics} Raising $\mu$ steepens the line. For $t>0$ the accepted set $\{g\le kt\}$ grows: added points have $t>0$ (raise $\tau$) and $g\in(k_{\text{old}}t,k_{\text{new}}t]$, $g>0$ (raise $\gamma$). For $t<0$ it shrinks: removed points have $t<0$ (removing them raises $\tau$) and $g\le kt<0$ (removing negative costs raises $\gamma$). So $\tau(\mu)$, $\gamma(\mu)$ increase. Lower $A$ requires higher pledgeable income $[p_H+\tau](R-B/\dpp)$, hence higher $\tau$, hence higher $\mu$: $\tau$ and $\gamma$ rise as $A$ falls.

\step{4}{(iii) Implementation} $x=1$, $t>0$: give the right to investors (they gain, bear no cost). $x=1$, $t<0$ (so $g\le kt<0$, a private benefit): give it to the entrepreneur, who takes it since $|g|\ge k|t|>|t|R\ge|t|R_b$. $x=0$, $t>0$ ($g>kt>tR_b$): entrepreneur, who declines. $x=0$, $t<0$: investors, who decline. No renegotiation: where the rule departs from first best ($x=1$, $t>0$, $tR<g\le kt$; or $x=0$, $t<0$, $kt<g\le tR$) the decision sits with investors and the entrepreneur, having invested all her cash, cannot compensate them (footnote~17, \S10.2.2).

\step{5}{(iv) Common $g>0$} The rule \eqref{r8e14-rule} becomes $t\ge t^*\equiv g/k(\mu)>0$, so
\begin{align}
  \gamma &= g\,[1-H(t^*)], \qquad \tau=\int_{t^*}^{\infty}t\,dH(t). \nonumber
\end{align}
Differentiating in $t^*$: $d\gamma/dt^*=-g\,h(t^*)$ and $d\tau/dt^*=-t^*h(t^*)$. Hence
\begin{align}
  \frac{d\gamma}{d\tau} &= \frac{-g\,h(t^*)}{-t^*h(t^*)}=\frac{g}{t^*}, \label{r8e14-d1}\\
  \frac{d^2\gamma}{d\tau^2} &= \frac{d}{dt^*}\Big(\frac{g}{t^*}\Big)\cdot\frac{dt^*}{d\tau} = \Big(-\frac{g}{t^{*2}}\Big)\Big(-\frac{1}{t^*h(t^*)}\Big)=\frac{g}{t^{*3}h(t^*)}>0. \label{r8e14-d2}
\end{align}

\step{6}{Numbers} $g=1$, $t\sim U[-0.1,0.3]$ ($h=2.5$), $R=10$, $B/\dpp=4$, $\mu=0.5$: $k=10+0.5(6)=13$, $t^*=1/13=0.0769$, $\gamma=2.5(0.3-1/13)=0.5577$, $\tau=1.25(0.09-1/169)=0.1051$; $d\gamma/d\tau=13$, $d^2\gamma/d\tau^2=1/(t^{*3}\cdot2.5)=878.8$ (finite differences agree). At $\mu=1$, $k=16$: $\tau$ and $\gamma$ both higher.

\answer{(i)} $x(t,g)=1$ iff $tR-g+\mu t(R-B/\dpp)\ge0$, i.e.\ below the line $g=[R+\mu(R-B/\dpp)]t$; $R_b=B/\dpp$.
\answer{(ii)} Lower $A$ $\Rightarrow$ higher $\mu$ $\Rightarrow$ steeper line $\Rightarrow$ $\tau$ and $\gamma$ both increase.
\answer{(iii)} $t>0$ rights to investors if $x=1$, to the entrepreneur if $x=0$; $t<0$ rights to the entrepreneur if $x=1$, to investors if $x=0$; not renegotiated.
\answer{(iv)} $d\gamma/d\tau=g/t^*$ and $d^2\gamma/d\tau^2=g/[t^{*3}h(t^*)]>0$.

\intuition Control rights are sold to investors in order of ``value to investors per unit of pain to the entrepreneur''. The cheapest ones go first, so each extra unit of pledgeable probability costs the entrepreneur more private cost: a convex cost of raising funds through control.

\pitfalls
\begin{itemize}
  \item Forgetting the $t<0$ half-plane: the financing need also makes her give up some value-reducing perks she would otherwise take.
  \item Treating $x$ as continuous in $\mu$ without noting that the line passes through the origin (no intercept).
\end{itemize}
\bookans Matches the book's answer (pp.~614--615): same cutoff line, rotation as $A$ falls, implementation, and $d\gamma/d\tau=g/t^*$ with $d^2\gamma/d\tau^2>0$.
\refbib{Tirole \S10.2.2 ``Multiple Control Rights'', pp.~392--393 (footnotes 15, 17); Answers, pp.~614--615.}

% =============================================================================
\exsec{15. Allocation of control over liquidation --- Tirole 10.2, pp.~418--419}
%% BEGIN-STATEMENT 15
\stmtbegin
\textbf{Source:} Tirole, Ch. 10, Exercise 10.2, pp.~418--419 (PDF 431) · Answer in book: yes,
pp.~611--612

\textbf{Exercise 10.2 (allocation of control and liquidation policy).} This exercise considers the
allocation of a control right over liquidation. As described in Figure 10.9, the framework has three
dates: date 0 (financing and investment), date 1 (choice of liquidation), and date 2 (payoff in the
case of continuation). There is moral hazard in the case of continuation. As usual, there is
universal risk neutrality, the entrepreneur is protected by limited liability, and the investors
demand a rate of return equal to 0.

Figure 10.9 (timeline): \emph{Date 0} --- financing (investors contribute \(I - A\)). → \emph{Date
1} --- either liquidation (verifiable liquidation value \(L\); entrepreneur obtains
\(U_b \geqslant 0\) in alternative job) or continuation. → After continuation: moral hazard (choice
of \(p \in \{p_L, p_H\}\)). → \emph{Date 2} --- profit (\(R\) with probability \(p\), 0 with
probability \(1 - p\)).

One will assume that the variables \((p_L, p_H, R, B)\) in the case of continuation are known
\emph{ex ante}. As usual, misbehaving (choosing probability \(p_L\)) yields a private benefit
\(B > 0\) to the entrepreneur. Let

\[\rho_0 \equiv p_H\left(R - \frac{B}{\Delta p}\right)\]

and

\[\rho_1 \equiv p_H R.\]

In contrast, the liquidation proceeds \(L\) and the fallback option \(U_b^0\) for the entrepreneur
may be \emph{ex ante} random, even though they become common knowledge at date 1 before the
liquidation decision. Lastly, \(L\) is fully pledgeable to investors while none of \(U_b^0\) is.

\begin{enumerate}
\def\labelenumi{(\roman{enumi})}
\item
  Solve for the optimal complete (state-contingent) contract, assuming that a court is able to
  directly verify \(\omega \equiv (L, U_b^0)\) (and the profit in the case of success) and to
  enforce the contract specifying the probability of continuation \(x(\omega) \in [0, 1]\) and the
  allocation of \(L\) and \(R\) between the investors and the entrepreneur.
\item
  Assume from now on that,
\end{enumerate}

\[\text{for all } \omega, \quad U_b^0 \leqslant \rho_1 - \rho_0.\]

That is, in the absence of a ``golden parachute'' given to the entrepreneur in the case of
liquidation, the entrepreneur always prefers to continue. Compare the sets \(\Omega^{FB}\) and
\(\Omega^{SB}\) of states of nature in which continuation is optimal in the absence and presence of
financing constraint. How does \(\Omega^{SB}\) vary with the entrepreneur's net worth \(A\)? (A
diagram will help.)

\begin{enumerate}
\def\labelenumi{(\roman{enumi})}
\setcounter{enumi}{2}
\tightlist
\item
  From now on, assume that the court observes neither \(L\) nor \(U_b^0\). Only the entrepreneur and
  the investors do. The remaining questions look at how far one can go toward the implementation of
  the optimal full-observability contract described in (i) using a simple allocation of the control
  right concerning liquidation.
\end{enumerate}

One will focus on the case in which \(\Omega^{SB}\) (see question (ii)) is strictly included in
\(\Omega^{FB}\), and so inefficient liquidation is required.

Suppose first that the entrepreneur has the control right and that renegotiation occurs once
\(\omega\) is realized. Argue that

\[\Omega^{EN} = \Omega^{FB},\]

where \(\Omega^{EN}\) is the set of states of nature over which continuation occurs under
entrepreneur control. Conclude that the project is then not financed.

\begin{enumerate}
\def\labelenumi{(\roman{enumi})}
\setcounter{enumi}{3}
\tightlist
\item
  \emph{Investor control.} Perform the analysis of question (iii) in the case of investor control in
  the absence of a golden parachute (the initial contract does not provide for any compensation for
  the entrepreneur in the case of liquidation). Suppose that the entrepreneur does not keep any
  savings. Show that
\end{enumerate}

\[\Omega^{IN} \subset \Omega^{SB},\]

where \(\Omega^{IN}\) is the set of states of nature over which continuation occurs under investor
control. Is the project financed?

\begin{enumerate}
\def\labelenumi{(\alph{enumi})}
\setcounter{enumi}{21}
\tightlist
\item
  \emph{Investor control with golden parachute.} Argue that a positive golden parachute (\(r_b > 0\)
  given to the entrepreneur in the case of liquidation) is optimal when investors have control.
\end{enumerate}

\paragraph{Context needed to solve it}

\emph{From Tirole §3.2.1--3.2.2 ``The Fixed-Investment Model'' and ``The Lenders' Credit Analysis'',
pp.~115--117 (PDF 128--130) --- the ``as usual'' moral-hazard stage:}

\emph{Project.} If undertaken, the project either succeeds, that is, yields verifiable income
\(R > 0\), or fails, that is, yields no income. The probability of success is denoted by \(p\). The
project is subject to moral hazard. The entrepreneur can ``behave'' (``work,'' ``exert effort,''
``take no private benefit'') or ``misbehave'' (``shirk,'' ``take a private benefit'') {[}\ldots{]}.

Behaving yields probability \(p = p_H\) of success and no private benefit to the entrepreneur, and
misbehaving results in probability \(p = p_L < p_H\) of success and private benefit \(B > 0\)
(measured in units of account) to the entrepreneur. {[}\ldots{]} Let \(\Delta p \equiv p_H - p_L\).

\emph{Preferences and the loan agreement.} Both the borrower and the potential lenders (or
``investors'') are risk neutral. For notational simplicity, there is no time preference; the rate of
return expected by investors (which is also the riskless rate, due to risk neutrality) is taken to
be 0. The borrower is protected by limited liability, and so her income cannot take negative values.

Lenders behave competitively in the sense that the loan, if any, makes zero profit. {[}\ldots{]} In
the case of success, the two parties share the profit \(R\); \(R_b\) goes to the borrower and
\(R_l\) to the lenders. {[}\ldots{]}

{[}\ldots{]} the borrower will therefore behave if the following ``incentive compatibility
constraint'' is satisfied:

\[p_H R_b \geqslant p_L R_b + B \quad\text{or}\quad (\Delta p)R_b \geqslant B. \qquad (\text{IC}_b)\]

From this incentive compatibility constraint we infer that the highest income in the case of success
that can be pledged to the lenders without jeopardizing the borrower's incentives is

\[R - \frac{B}{\Delta p}.\]

The \emph{(expected) pledgeable income} is then

\[\mathcal{P} = p_H\left(R - \frac{B}{\Delta p}\right).\]
\stmtend
%% END-STATEMENT

\asked (i) The optimal state-contingent contract when $\omega=(L,U_b^0)$ is verifiable. (ii) Compare the first-best and second-best continuation sets; how $\Omega^{SB}$ moves with $A$. (iii) Entrepreneur control with renegotiation: show $\Omega^{EN}=\Omega^{FB}$ and that the project is not financed. (iv) Investor control without golden parachute: show $\Omega^{IN}\subset\Omega^{SB}$; is it financed? (v) Why a positive golden parachute is optimal under investor control.

\approach A complete-contract benchmark first (an LP in $x(\omega)$ with a shadow price on funds), then the incomplete-contract question of Aghion--Bolton: which allocation of a single control right, plus ex post renegotiation, comes closest. The key asymmetry: investors can pay the entrepreneur (they have cash); the entrepreneur, having invested everything, cannot pay investors --- unless a golden parachute gives her something to give up.

\notation
\begin{itemize}
  \item $x(\omega)\in[0,1]$: continuation probability; $\rho_1=p_HR$: continuation value; $\rho_0=p_H(R-B/\dpp)$: pledgeable continuation value; $\rho_1-\rho_0=p_HB/\dpp$: entrepreneur's rent.
  \item $\nu\ge0$: multiplier on breakeven, $\mu\equiv1+\nu\ge1$; $r_b$: golden parachute.
\end{itemize}

\step{1}{(i) Contract terms} Pay the entrepreneur $0$ on failure and give investors all of $L$: both raise pledgeable income without affecting incentives. Set $R_b=B/\dpp$ when financing binds. Then the entrepreneur gets $\rho_1-\rho_0$ on continuation and $U_b^0$ on liquidation; investors get $\rho_0$ or $L$.

\step{2}{(i) Program} Investors break even, so $U_b$ equals total value minus $I$:
\begin{align}
  \max_{x(\cdot)}\ & E_\omega\big[x\rho_1+(1-x)(L+U_b^0)\big]-I \quad\text{s.t.}\quad E_\omega\big[x\rho_0+(1-x)L\big]\ge I-A. \nonumber
\end{align}
Pointwise derivative of the Lagrangian with respect to $x(\omega)$:
\begin{align}
  \rho_1-L-U_b^0+\nu(\rho_0-L). \nonumber
\end{align}
Continue iff it is $\ge0$; rearranging with $\mu=1+\nu$:
\begin{align}
  x^{SB}(\omega)=1 &\iff \rho_1-U_b^0\ge\mu L-(\mu-1)\rho_0. \label{r8e15-sb}
\end{align}

\step{3}{(ii) First best vs.\ second best} First best is $\mu=1$: continue iff $L+U_b^0\le\rho_1$. In the $(L,U_b^0)$ plane the two boundaries are $U_b^0=\rho_1-L$ and $U_b^0=\rho_1+(\mu-1)\rho_0-\mu L$; both pass through $(\rho_0,\rho_1-\rho_0)$ and differ by $(\mu-1)(\rho_0-L)$. For $L\le\rho_0$, under assumption (ii), both say continue: $\mu L-(\mu-1)\rho_0=\rho_0+\mu(L-\rho_0)\le\rho_0\le\rho_1-U_b^0$. For $L>\rho_0$ the SB boundary is lower. Hence $\Omega^{SB}\subseteq\Omega^{FB}$, the difference being states with $L>\rho_0$ and $L\le\rho_1-U_b^0<\mu L-(\mu-1)\rho_0$ (Figure). As $A$ falls the required investor income rises; income $E[\rho_0+(L-\rho_0)\mathbf 1\{\text{liquidate}\}]$ grows only by liquidating more $L>\rho_0$ states, i.e.\ by raising $\mu$: the SB line rotates clockwise around the pivot and $\Omega^{SB}$ shrinks.

\begin{center}
\begin{tikzpicture}[scale=1.15]
  \draw[->] (0,0) -- (4.4,0) node[right] {$L$};
  \draw[->] (0,0) -- (0,2.4) node[above] {$U_b^0$};
  \draw[dashed] (0,1.6) -- (4.2,1.6) node[right] {\footnotesize $\rho_1-\rho_0$};
  \draw[thick] (1.6,2.4) node[above] {\footnotesize FB: $L+U_b^0=\rho_1$} -- (4,0);
  \draw[thick,blue!60!black] (2.0,2.2) node[right] {\footnotesize SB (slope $-\mu$)} -- (3.47,0);
  \fill (2.4,1.6) circle (1.2pt) node[above right] {\footnotesize $(\rho_0,\rho_1-\rho_0)$};
  \draw[dotted] (2.4,0) node[below] {\footnotesize $\rho_0$} -- (2.4,1.6);
  \node at (1.0,0.7) {\footnotesize continue};
  \node at (3.4,1.35) {\footnotesize liquidate};
  \node[blue!60!black,right] at (4.1,0.9) {\footnotesize $\Omega^{FB}\setminus\Omega^{SB}$};
  \draw[->,blue!60!black] (4.15,0.85) -- (3.35,0.4);
\end{tikzpicture}
\end{center}

\step{4}{(ii) Example} $\rho_1=4$, $\rho_0=2.4$ ($p_H=0.8$, $\dpp=0.3$, $B=0.6$, $R=5$), $U_b^0=1$, $L\sim U[0,4]$, $I=4$. First best: continue iff $L\le3$. Investor income with cutoff $\ell\ge\rho_0$: $P(\ell)=\tfrac{\ell}{4}(2.4)+\int_\ell^4\tfrac{L}{4}dL=0.6\ell+(16-\ell^2)/8$; $P(3)=2.675$. With $A=1.3$ ($I-A=2.7$): $0.6\ell+(16-\ell^2)/8=2.7\iff\ell^2-4.8\ell+5.6=0\iff\ell=2.8$ (the root $\ell\ge\rho_0$). From \eqref{r8e15-sb} at the cutoff, $3=2.8\mu-2.4(\mu-1)$, so $\mu=1.5$. With $I-A=2.71$: $\ell=2.6828$, $\mu=2.1213$ --- $\Omega^{SB}$ shrinks.

\step{5}{(iii) Entrepreneur control} Without a parachute she always wants to continue (she gets $\rho_1-\rho_0\ge U_b^0$). Renegotiation is efficient ex post because investors have cash: where $L+U_b^0>\rho_1$, investors pay her at least $\rho_1-\rho_0-U_b^0$ to liquidate, gaining $L-\rho_0$, and $L-\rho_0>\rho_1-\rho_0-U_b^0$; where continuation is efficient there are no gains from trade. So $\Omega^{EN}=\Omega^{FB}$. Investors then get at most $\rho_0$ on $\Omega^{FB}$ and at most $L-(\rho_1-\rho_0-U_b^0)\le L$ outside. Compared with the SB contract, whose income is exactly $I-A$, they lose at least $E[(L-\rho_0)\mathbf 1\{\Omega^{FB}\setminus\Omega^{SB}\}]>0$: not financed. Example: $0.75(2.4)+\tfrac14\int_3^4(L-0.6)dL=1.8+0.725=2.525<2.7$.

\step{6}{(iv) Investor control, no parachute} Investors get $p_H(R-R_b)\le\rho_0$ from continuing and $L$ from liquidating; the entrepreneur has no cash to bribe them, and can at most cut $R_b$ to $B/\dpp$. So they continue iff $L\le\rho_0$: $\Omega^{IN}=\{L\le\rho_0\}$. By Step 3 every such state is in $\Omega^{SB}$, and $\Omega^{SB}$ also contains states with $L$ slightly above $\rho_0$: $\Omega^{IN}\subset\Omega^{SB}$. Investors get the SB income plus $L-\rho_0>0$ on $\Omega^{SB}\setminus\Omega^{IN}$: financed, but with excessive liquidation. Example: $P(2.4)=2.72\ge2.7$; NPV $0.08$ vs.\ $0.12$ (SB) and $0.125$ (FB).

\step{7}{(v) Golden parachute} With $r_b>0$ paid out of $L$ on liquidation, investors liquidate iff $L-r_b>\rho_0$. If continuation is efficient but $L-\rho_0>r_b$, she can offer at most $r_b$: inefficient liquidation remains. If $L-\rho_0<r_b$ but liquidation is efficient, $(\rho_1-\rho_0)-U_b^0\le L-\rho_0$, she gives up part of $r_b$ and the parties renegotiate to liquidation. A small $r_b$ thus moves states with $L\in(\rho_0,\rho_0+r_b]$, all efficient continuations, from liquidation to continuation: NPV rises. Investor income falls continuously in $r_b$, and in (iv) it had slack, so a small $r_b>0$ keeps financing. Example: $r_b=0.05$, cutoff $2.45$, income $P(2.45)-0.05\cdot\tfrac{4-2.45}{4}=2.7197-0.0194=2.7003\ge2.7$, NPV $0.0872>0.08$; $r_b=0.10$ would give $2.6813<2.7$.

\answer{(i)} $R_b=B/\dpp$, $0$ on failure, all of $L$ to investors; continue iff $\rho_1-U_b^0\ge\mu L-(\mu-1)\rho_0$, $\mu\ge1$.
\answer{(ii)} $\Omega^{SB}\subseteq\Omega^{FB}$ (equal for $L\le\rho_0$, smaller for $L>\rho_0$); $\Omega^{SB}$ shrinks as $A$ falls ($\mu$ rises).
\answer{(iii)} $\Omega^{EN}=\Omega^{FB}$; investors get less than $I-A$: no financing.
\answer{(iv)} $\Omega^{IN}=\{L\le\rho_0\}\subset\Omega^{SB}$; financed, with too much liquidation.
\answer{(v)} A small $r_b>0$ restores efficient continuation for $L$ just above $\rho_0$ while keeping investors at breakeven: optimal.

\intuition Whoever holds control, the other side bribes efficiently only if it has a currency. Investors always do, so entrepreneur control reproduces the first best --- which promises too little to investors. The entrepreneur has no currency, so investor control liquidates whenever $L>\rho_0$ --- too often. A golden parachute is the currency that lets her ``buy'' continuation, a reason for parachutes beyond managerial entrenchment.

\pitfalls
\begin{itemize}
  \item Comparing $L$ with $\rho_1$ for investors: they only see $\rho_0$, the pledgeable part.
  \item Forgetting renegotiation under entrepreneur control: it delivers $\Omega^{FB}$, not ``always continue''.
\end{itemize}
\bookans Matches the book's answer (pp.~611--612): rule $\rho_1-U_b^0\ge-(\mu-1)\rho_0+\mu L$; $\Omega^{SB}\subset\Omega^{FB}$, shrinking with $A$; $\Omega^{EN}=\Omega^{FB}$ and no financing; $x^{IN}=1$ iff $\rho_0\ge L$, financed; renegotiation conditions for $r_b$. The numerical example is added here.
\refbib{Tirole \S10.2.1, pp.~389--391; \S3.2, pp.~115--117; Answers, pp.~611--612.}

\end{document}
